\documentclass[UTF8,a4paper,11pt,fontset=none]{ctexart}
\usepackage{amsmath,amssymb,mathtools,bm}
\usepackage{fontspec}
\usepackage{geometry}
\usepackage{array,booktabs,longtable,tabularx}
\usepackage{enumitem}
\usepackage{xcolor}
\usepackage{fancyhdr}
\usepackage{hyperref}
\usepackage{xurl}
\geometry{top=24mm,bottom=24mm,left=23mm,right=23mm,headheight=15pt}
\setmainfont{Latin Modern Roman}
\setsansfont{Latin Modern Sans}
\setmonofont{Latin Modern Mono}
\setCJKmainfont[Path={latex_assets/fonts/},Extension=.otf,UprightFont=*-Regular,BoldFont=*-Bold,ItalicFont=*-Regular,BoldItalicFont=*-Bold]{SourceHanSerifSC}
\setCJKsansfont[Path={latex_assets/fonts/},Extension=.otf,UprightFont=*-Regular,BoldFont=*-Bold,ItalicFont=*-Regular,BoldItalicFont=*-Bold]{SourceHanSansSC}
\XeTeXlinebreaklocale "zh"
\XeTeXlinebreakskip=0pt plus 1pt
\setlength{\parindent}{2em}
\setlength{\parskip}{0.35em}
\setlength{\emergencystretch}{3em}
\linespread{1.28}
\setlist{nosep}
\hypersetup{unicode=true,colorlinks=true,linkcolor=blue!45!black,urlcolor=blue!55!black,pdftitle={张靖皋猫道十四项低阶模态独立理论建模与盲复现方案}}
\pagestyle{fancy}
\fancyhf{}
\fancyhead[L]{张靖皋猫道理论建模与验证}
\fancyhead[R]{独立理论验证}
\fancyfoot[C]{\thepage}
\renewcommand{\headrulewidth}{0.4pt}
\begin{document}
\hypersetup{pageanchor=false}
\begin{titlepage}
\centering
\vspace*{0.16\textheight}
{\fontsize{24pt}{34pt}\selectfont\bfseries 张靖皋猫道十四项低阶模态独立理论建模与盲复现方案\par}
\vspace{1.4em}
{\fontsize{14pt}{22pt}\selectfont ——以扭转族为主线，基于图纸转录、现有实体报告、公开材料标准与独立复算，不调用B00及既有MCT数值\par}
\vfill
\begin{minipage}{0.88\textwidth}
\centering
版本：v1.2（独立理论验证与证据审计定稿版）\par
日期：2026-08-08\par
验证状态：L0解析解、离散收敛和机制单元验证完成；在不使用B00、既有MCT数值或目标频率反标参数的条件下，已独立覆盖4个横弯、4个竖弯和3个SIDE家族，共11/14个家族，组合平均绝对相对误差约2.51\%；L1/L2分段连续体尚未完成组装，因此三项扭转及严格“前14阶同序复现”仍未签发\par
\end{minipage}
\vspace*{0.10\textheight}
\end{titlepage}
\pagenumbering{Roman}
\tableofcontents
\clearpage
\pagenumbering{arabic}
\setcounter{page}{1}
\hypersetup{pageanchor=true}
\phantomsection
\section*{0. 本文要解决的问题、硬约束与结论先行}
\addcontentsline{toc}{section}{0. 本文要解决的问题、硬约束与结论先行}
本文把“复现附件2-3的十四项低阶物理模态”解释为：在不使用B00、既有MCT节点/单元/质量/刚度/索力/频率、既有降阶矩阵及任何由既有模态反算出的参数的前提下，依据原稿对175页《图纸汇总-2026.08.05.pdf》的转录，并以本轮可直接核查的《猫道图纸1225》《猫道结构复核计算报告0324》和《猫道结构抗风性能试验报告》交叉验证，重新建立一个具有明确物理来源的四跨双幅猫道模型。175页新图纸原文件本轮未随附件或归档出现，故凡只见于原稿转录而未被现有实体文件交叉闭合的数据，均保留版本警示，绝不把旧图默认成新图。盲验协议拆成三层：TARGET-FAMILY只保存十四个物理家族名称，可在冻结前用于规定输出分类；TARGET-SHAPE保存附件振型图或振型向量，只能在候选谐波和特征根全部计算并冻结后用于形态分类，不得用于选基、择阶或反标参数；TARGET-FREQ只保存数值频率，必须在几何、质量、材料、边界、静力平衡、离散和降阶参数全部冻结后才加载。因用户提供的原稿已经明列目标频率和振型图，本次执行属于“计算输入隔离的后验理论复核”，不能冒充严格双盲试验；未来正式运行应把TARGET-SHAPE与TARGET-FREQ置于求解器不可读的独立文件并在冻结前保存哈希。本文不把“双MCT”理解为复用两个已有MCT模型，而把它保留为一种拓扑思想：上、下游两幅猫道分别建成独立子结构，二者之间仅通过离散横向通道及真实接口耦合，任何对称关系都由计算结果呈现，不在建模阶段强制施加。\par
基于现有材料追索后的事实，最值得优先检验的扭转机制有六个。第一，单幅猫道由16根面层承重索、门架承重索、扶手索、横梁、面层、71道高门架及其底梁共同组成，扭转惯量和扭转刚度具有显著的竖向与横向偏心；第二，两幅猫道中心线间距42.90 m，横向通道属于离散重构件，南侧H1—H12详图支持约49.655 m总长、约1.50 m宽和约1.70 m高的三角桁架构形，不能压缩成只保留一个平动方向的单弹簧；北侧九道的位置已由总体尺寸链全部恢复，是否与南侧完全同型仍需BOM确认；第三，H1—H12通道局部倾角从0°变化到正18.89°并跨塔变号至−20.24°，说明通道局部坐标、纵向—竖向耦合和站位差异必须逐站保留；第四，四跨长度为660 m、2300 m、717 m、503 m，南侧另有辅助塔，整桥纵向不存在全局镜像对称；第五，门架存在M1/M2类型差异且间距按51 m、57 m、60 m等多种节距布置，塔区、跨中和边跨的局部刚度、质量与连接条件均不周期；第六，每幅16根面层承重索中有1根智慧芯绳，其设计截面积比普通绳小7.71\%，其实际横向安装位置会形成真实的单索材料不对称，是近重根旋转与微小分裂的高价值探针。上述机制构成待证伪假设，其重要性不能仅凭秩一杆件的局部耦合比例预先下结论；门架和通道的多端口刚度、偏心质量、局部畸变、索力/表观轴向刚度及边界均须通过消融和独立数据比较。本轮已从施工图、复核报告和抗风报告找回原稿误列缺项：成型猫道主跨垂度227.300 m、全线21道通道位置、单幅71道门架位置、设计用钢丝绳\(E=1.20\times10^5\ \mathrm{MPa}\)\allowbreak{}与截面积、各跨曲线长与无应力长度，以及恒载索力审计值。仍需证据化的事项已经转化为具体行动：取得产品证书或样段试验确认工作区间切线\(EA_t\)\allowbreak{}，从构造/局部模型或试验生成夹具、网片、门架和通道客观多端口矩阵，核实北侧九道BOM与姿态及智慧芯安装位置，确认鞍座粘滑状态，并盘点模态测试时附加质量。本文已经完成L0与机制单元验证，但尚未组装L1/L2分段连续体，故不发布伪精确的十四项数字，也不使用TARGET-FREQ反标缺项。\par
证据等级提示：上段H1—H12的模块尺寸、单重和逐站倾角来自原稿对175页新图的转录，对应原PDF本轮未取得，因此只能称“转录详图所示”，证据状态为TRANSCRIBED-DRW-A-UNVERIFIED。本轮直接可核的一手证据是2025实体旧图、2026复核报告和抗风报告；转录数值只在带警示的诊断分支中使用。\par
\phantomsection
\section*{1. 数据隔离、来源等级与防止目标泄漏的规则}
\addcontentsline{toc}{section}{1. 数据隔离、来源等级与防止目标泄漏的规则}
所有数据按“五类输入证据加一类目标”管理，其中图纸证据再分A/B两级。DRW-A表示图纸直接标注的尺寸、数量、型号、单重、总重和布置；DRW-B表示由图纸尺寸链直接推导且可复核的量，例如16根面层承重索的横向坐标、二次矩及全线通道坐标；CALC-INPUT表示复核报告中作为计算前提明列的控制点、材料或荷载输入，必须与该报告的计算输出分开；STD表示公开标准给出的钢材常数或产品验收要求；TEST表示必须由材料证书、张拉记录、索力测试、静载试验、局部构件试验或实测几何获得的量；ASSUMP表示尚无独立证据时建立的对照假设，只允许形成敏感性区间。复核报告中的成型曲线、索力、无应力长度、位移和反力另标CALC-AUDIT，只在独立模型冻结后用于差异诊断，不回写输入。TARGET-FAMILY表示冻结前可知的物理家族标签，TARGET-FREQ表示冻结后才加载的数值频率。任何模型输入都必须附带source\_tag、drawing\_no/page、unit、value、uncertainty、model\_mapping和是否允许更新等字段；TARGET-FREQ文件与模型输入文件分开存放，冻结时分别计算SHA-256，求解器阶段不得读取TARGET-FREQ文件。\par
补充定义：TRANSCRIBED-DRW-A-UNVERIFIED表示只存在于原稿转录、但对应原PDF本轮未取得的图纸直读量；该类量只能进入带警示的诊断分支，不得视为已核一手证据。TARGET-SHAPE表示振型图或振型向量；它与TARGET-FREQ一样与模型输入隔离，只能在候选解全部计算并冻结后作形态分类，不能参与参数识别。\par
原稿登记的新图文件为《图纸汇总-2026.08.05.pdf》，共175页、A3页面，并记录下列SHA-256。该原PDF本轮未随附件或项目归档出现，因此文件名、页码、图号和哈希均属原稿转录元数据，状态为TRANSCRIBED-DRW-A-UNVERIFIED，不得冒充本轮已核源文件： \nolinkurl{5b6231a766c549db07539df85d8c14b99b816dc0772abb1ab19d7deb782b0bd9}\par
明确排除的数值来源包括：B00模型的任意几何、材料、截面、质量、约束、索力、模态和振型；既有MCT模型及其节点、单元、属性、集中质量、索力、刚度矩阵和凝聚矩阵；既有降阶模型为贴合频率而采用的校准刚度或校准质量；从附件2-3频率反算得到的索力、EA、转动弹簧或横向通道等效刚度。可以使用的理论来源包括非线性索单元、预应力模态、Timoshenko/Vlasov梁、静力凝聚与Craig—Bampton部件模态综合等公开理论；可以使用的工程数据包括本次图纸、正式材料标准、制造商材质证明和独立试验。\par
\phantomsection
\section*{2. 从图纸转录、实体旧图与复核报告提取的模型事实}
\addcontentsline{toc}{section}{2. 从图纸转录、实体旧图与复核报告提取的模型事实}
\phantomsection
\subsection*{2.1. 总体几何与非对称布置}
\addcontentsline{toc}{subsection}{2.1. 总体几何与非对称布置}
节距口径预先更正：最南503 m跨依实体图示\(5200+7\times5700+5200=50300\)\allowbreak{}的闭合尺寸链采用52 m端距。本节原始摘要出现的52.5 m只是OCR说明文字冲突，不是模型采用值；详见2.7节逐站闭合。\par
图纸M00-01（PDF第1页）给出四跨连续猫道：北侧660 m段、主跨2300 m、南侧717 m段、最南503 m段，总长4180 m；两幅猫道中心线间距42.90 m。传统跨名存在版本冲突：2025施工图MD0-01和原稿所转录的新图把717 m称“南辅跨”、503 m称“南边跨”，2026复核报告第4页文字则把717 m称“南边跨”、503 m称“南辅跨”。两版长度、端点和尺寸链一致，故本文和机器表只以“南侧717 m段”“最南503 m段”为主键，传统名称仅作来源别名，不据此改变位置。原稿把总体立面标出的255.56 m直接当作猫道主跨垂度；现有材料追溯表明该尺寸属于主缆/桥梁控制线。复核报告表1-5（PDF第9页）把猫道承重索成型线形作为控制点找形，塔顶控制高程为340.600 m、主跨跨中控制高程为113.300 m；表1-9的逐点成型坐标在PDF第23页给出北塔K17+542.679、340.600 m，在第26页给出主跨跨中K18+686.000、113.300 m，在第22页给出南塔K19+829.321、340.600 m，因此成型猫道主跨垂度为227.300 m。表1-5把主跨跨中桩号印成K19+836.000；报告表1-6第10页把K19+836.000正确用于南塔门架分跨点，表1-9又以逐点序列闭合K18+686.000跨中，故将表1-5桩号登记为排版笔误，不影响两处一致的113.300 m控制高程。猫道中心线与主缆中心线近似平行，主跨竖向间距约1.7 m，边跨约1.3—1.8 m。横向通道总数为21道，其中北侧660 m段3道、主跨13道、南侧717 m段3道、最南503 m段2道；主跨站距为(178+3×171+6×153+3×171+178)m，北侧660 m段为(159+171+171+159)m，南侧717 m段为(178.5+180+180+178.5)m，最南503 m段为(166+171+166)m。门架总数为单幅71道：主跨41道、北侧660 m段11道、南侧717 m段11道、最南503 m段8道；主跨跨中附近约51 m一榀，两侧约57 m一榀，塔旁首间约64 m，其他跨还出现45 m、52 m、58.5 m、60 m等节距。52.5 m仅保留为最南503 m段说明文字的OCR冲突，不进入站位。以上事实直接否定整桥纵向半模型和统一周期单元。\par
名义尺寸链与自由索段必须分开。总体图的名义北塔和南塔位置为\(x=660\)\allowbreak{} m与\(x=2960\)\allowbreak{} m，对应名义主跨2300 m；复核报告表1-5/1-9给出的主跨侧实际鞍点为K17+542.679和K19+829.321，即\(x=666.679\)\allowbreak{} m与\(x=2953.321\)\allowbreak{} m，自由索段长度为2286.642 m。主跨跨中H1相对名义北塔为1150 m，相对实际北侧鞍点为1143.321 m。5.1节的等高抛物线自重频率中跨长会代数消去，故名义/自由段差异不改变该L0尺度；成型线拟合、精确悬链线、站位形函数和分段传递矩阵必须并列记录两种坐标口径，不能混用。\par
若把两个主跨下拉点而不是鞍点视为振动自由段边界，表1-5给K17+553.300和K19+818.700、高程334.645 m与334.537 m，对应水平长2265.400 m、相对平均端高的跨中矢高约221.291 m。这是第三种有构造含义的边界定义，不应与2286.642 m鞍到鞍段静默混用；它进入B1/B2/B3边界敏感性，由下拉连接的实际约束状态筛选，而不是由目标频率择优。\par
\phantomsection
\subsection*{2.2. 单幅猫道连续索质量}
\addcontentsline{toc}{subsection}{2.2. 单幅猫道连续索质量}
M00-01的钢丝绳材料表给出单幅猫道的四组连续索：16根φ50面层承重索，单位质量12.038 kg/m；2根φ36上扶手索，单位质量5.42 kg/m；4根φ20中、下扶手索，单位质量1.67 kg/m；6根φ50门架承重索，单位质量12.038 kg/m。由图纸直接得到的连续索线质量为\par
\[
\mu_{\mathrm{rope}}
=16\times 12.038+2\times 5.42+4\times 1.67+6\times 12.038
=282.356\ \mathrm{kg/m}.
\]
其中16根面层承重索的线质量为\par
\[
\mu_{\mathrm{floor}}=16\times 12.038=192.608\ \mathrm{kg/m}.
\]
版本审计发现，项目归档中的较早《猫道图纸1225》MD0-01（PDF第6页）把门架承重索列为6根\(\phi54\)\allowbreak{}、12.2 kg/m，因此该版四组连续索合计为283.328 kg/m；原稿所依据的2026.05版转录值把门架承重索列为6根\(\phi50\)\allowbreak{}、12.038 kg/m，复核报告第6页也采用\(\phi50\)\allowbreak{}门架承重索。主分支按时间更晚且被复核报告交叉支持的\(\phi50\)\allowbreak{}口径保留282.356 kg/m，旧版283.328 kg/m作为版本敏感性，不把两版绳径、质量和构造交叉拼接。两版连续索总质量相差0.972 kg/m（0.344\%）。这一数值只包含表中四组纵向钢丝绳，不包含横梁、踏步方木、钢丝网、门架、底梁、夹具、横向通道和施工附属物，因此不得直接当作全结构总线质量。\par
复核报告的荷载包是另一种口径：表1-1把底索单根质量列为11.989 kg/m，并把底部猫道系统恒载汇总为2.766 kN/m，按本文标准重力折合282.0535 kg/m；该汇总还包含扶手、网片、木条、小横梁和电缆等，表1-3又把六根门架承重索另列0.709 kN/m。故11.989与12.038 kg/m构成报告/施工图版本差异，2.766 kN/m也不能与上述“四组纯钢丝绳282.356 kg/m”直接相减或叠加。5.3节面内基线使用报告2.766 kN/m这一完整CALC-INPUT荷载包，2.2节纯钢丝绳算术只用于图纸质量审计。\par
\phantomsection
\subsection*{2.3. 面层承重索横向坐标与扭转二次矩}
\addcontentsline{toc}{subsection}{2.3. 面层承重索横向坐标与扭转二次矩}
M00-04（PDF第4页）给出单幅猫道总宽约5.78 m、面网大横梁有效宽约5.60 m，16根面层承重索分成左右各8根，组内间距0.260 m，内侧两组间净距1.700 m。以单幅猫道中心线为y=0，按图纸尺寸链得到名义横向坐标\par
\[
y_i\in\left\{\pm\left(0.85+0.26k\right)\,\middle|\,k=0,1,\ldots,7\right\}\ \mathrm{m}.
\]
故\par
\[
\sum_{i=1}^{16}y_i^2=55.24\ \mathrm{m^2}.
\]
仅计16根面层承重索的单位长度扭转质量二次矩下限为\par
\[
I_{\theta,\mathrm{floor}}
=\sum_{i=1}^{16}\mu_i y_i^2
=12.038\times55.24
=664.979\ \mathrm{kg\,m}.
\]
若16根索在某截面具有相同的有效水平张力分量H\_c，则其张力几何刚度对扭转的系数为\par
\[
C_{\theta,\mathrm{floor}}
=\sum_{i=1}^{16}H_c y_i^2
=55.24H_c\ \mathrm{N\,m^2}.
\]
这两个量均由图纸坐标与单位质量直接推导。上扶手索、中下扶手索和门架承重索具有更大的z偏心或不同的y位置，它们会进一步增加扭转惯量并引入平移—扭转耦合，最终坐标必须从门架详图/CAD中逐根提取。\par
\phantomsection
\subsection*{2.4. 横梁、踏步与门架离散质量}
\addcontentsline{toc}{subsection}{2.4. 横梁、踏步与门架离散质量}
M01-02（PDF第7页）给出单套面网大横梁质量116.93 kg，单幅共585套，总质量68.40405 t；M01-03（PDF第8页）给出单套面网小横梁质量30.18 kg，单幅共845套，总质量25.50210 t，并给出踏步方木50×30×2000、单重1.62 kg、单幅15000根，总质量24.30000 t。M01-05（PDF第10页）给出单套门架底梁质量306.98 kg，单幅71套，总质量21.79558 t。M04-02（PDF第142页）转录单套门架质量1123 kg且不含底梁，按71套算术上对应79.733 t；M04-03（PDF第143页）显示M1横梁为□160×160×6×7460、质量211.00 kg，M2横梁为□160×160×8×7460、质量275.18 kg，单件差64.18 kg。正式质量台账必须回到原始BOM核实1123 kg究竟已包含各类型横梁，还是仅为某一基准类型；若已包含，另加M1/M2差值会重复计量，若仅对应M1，则还需按M2数量增加\(64.18n_{\mathrm{M2}}\)\allowbreak{} kg。M1/M2具体站位必须按M04-01标签逐站建立。\par
较早《猫道图纸1225》的版本口径不同：MD4-01（PDF第91页）写单套门架829.3 kg且不含底梁，MD1-05（PDF第15页）写底梁313.46 kg，合计1142.76 kg/套；该PDF的MD4-01—MD4-07没有出现1123 kg、M1/M2或64.18 kg差值。主分支继续使用原稿所据2026.05版转录口径，旧版只形成完整的“829.3+313.46”对照，严禁取旧版829.3 kg再配新版306.98 kg底梁，或把新版M1/M2差值叠加到旧版门架。以上构件不应先摊成均匀线质量后永久使用，首轮解析模型可以用均匀化量做量级核对，正式模型必须把它们放回同一版本图纸的站位与实际偏心位置。\par
大横梁、小横梁、踏步、门架底梁和门架若临时除以4180 m，可得到16.365、6.101、5.813、5.214和19.075 kg/m的平均量级；这些平均值仅用于检查总质量，不允许替代正式离散布置。\par
\phantomsection
\subsection*{2.5. 横向通道几何、质量与局部角度}
\addcontentsline{toc}{subsection}{2.5. 横向通道几何、质量与局部角度}
M05-01—M05-12位于原稿所据175页新图PDF第151—162页。M05-01说明该标段覆盖主跨跨中及其以南的12道通道：主跨跨中至南主塔7道、南侧717 m跨3道、最南503 m跨2道。原稿转录的M05-02给出一整道通道总长49.655 m，两幅猫道中心线间距42.90 m，模块沿横桥向依次为7.435、9.540、9.540、6.165、9.540、7.435 m；通道桁架横断面约1.50 m宽、1.70 m高。原稿转录的M05-03给出9.54 m中间模块质量1100.9 kg且每道3套；M05-05给出7.435 m尾段模块质量865 kg且每道2套；M05-07给出6.165 m中间模块质量736.5 kg且每道1套。因此在不计护栏、端部接口、螺栓和局部附件时，新版转录所示单道通道的基本桁架模块质量为\par
\[
m_{\mathrm{passage,core}}
=3\times1100.9+1\times736.5+2\times865
=5769.2\ \mathrm{kg}.
\]
M05-10—M05-12另列护栏和钢丝网片等质量，正式质量台账必须逐项组装并检查“单个模块”“单个通道共若干套”等口径，避免把表内已汇总数量再次乘以套数。按5769.2 kg/道计算，H1—H12十二道详图范围内的核心模块合计69.2304 t。现有《猫道结构复核计算报告0324》表1-2把每幅猫道承担的半道横向通道设计恒载列为5065 kg，对应整道10130 kg；该值包含范围明显大于5769.2 kg核心模块，可作为CALC-INPUT质量上界和BOM核对对象，不能与核心模块质量叠加。北侧九道同级构件详图未出现在该标段中，但其位置已经由总体尺寸链完整恢复，不能再列为“位置未知”。以总体布置四跨北端K16+876.000为\(x=0\)\allowbreak{}并按北向南为正（该原点不是物理锚固点K16+852.105），位置主分支采用2025施工图/2026复核报告口径；同时先登记以下几何/BOM版本差异：\par
实体《猫道图纸1225》MD5-02给出的对应版本总长为49.920 m、控制/平面长度为49.720 m，模块单重也与上述新版转录不同；复核报告第111页采用约49.65 m、高1.75 m、宽1.5 m的倒三角桁架描述。三套数据属于版本差异，不能静默平均。本文的5769.2 kg只用于原稿新版转录的核心模块场景，10130 kg只用于复核报告整道设计恒载场景，实体旧图则单独作为几何/BOM版本审计；工程签发前必须取得175页新图原文件逐项复核。\par
全线21道主分支位置为：\par
北边跨3道：\(x=159,\allowbreak 330,\allowbreak 501\ \mathrm m\)\allowbreak{}；\par
主跨北半6道：\(x=838,\allowbreak 1009,\allowbreak 1180,\allowbreak 1351,\allowbreak 1504,\allowbreak 1657\ \mathrm m\)\allowbreak{}；\par
H1—H7（主跨跨中及以南）：\(x=1810,\allowbreak 1963,\allowbreak 2116,\allowbreak 2269,\allowbreak 2440,\allowbreak 2611,\allowbreak 2782\ \mathrm m\)\allowbreak{}，相对主跨北塔\(x_m=1150,\allowbreak 1303,\allowbreak 1456,\allowbreak 1609,\allowbreak 1780,\allowbreak 1951,\allowbreak 2122\ \mathrm m\)\allowbreak{}；\par
H8—H10（南侧717 m跨）：\(x=3138.5,\allowbreak 3318.5,\allowbreak 3498.5\ \mathrm m\)\allowbreak{}；\par
H11—H12（最南503 m跨）：\(x=3843,\allowbreak 4014\ \mathrm m\)\allowbreak{}。\par
下表把21道通道的位置和工程桩号逐项闭合。P01—P09是本文为北侧无详图标签的9道通道设置的审计号，H1—H12沿用现有详图标签。\par
\clearpage
\begingroup
\small
\setlength{\tabcolsep}{3pt}
\setlength{\LTleft}{\fill}
\setlength{\LTright}{\fill}
\begin{longtable}{|>{\raggedright\arraybackslash}p{0.0500\textwidth}|>{\raggedright\arraybackslash}p{0.1000\textwidth}|>{\raggedright\arraybackslash}p{0.1500\textwidth}|>{\raggedright\arraybackslash}p{0.1400\textwidth}|>{\raggedright\arraybackslash}p{0.1400\textwidth}|>{\raggedright\arraybackslash}p{0.2000\textwidth}|}
\hline
\textbf{序号} & \textbf{详图/审计号} & \textbf{跨段} & \textbf{跨内位置 (m)} & \textbf{总体坐标 \(x\)\allowbreak{} (m)} & \textbf{工程桩号} \\ \hline
\endfirsthead
\hline
\textbf{序号} & \textbf{详图/审计号} & \textbf{跨段} & \textbf{跨内位置 (m)} & \textbf{总体坐标 \(x\)\allowbreak{} (m)} & \textbf{工程桩号} \\ \hline
\endhead
1 & P01 & 北边跨 & 159 & 159 & K17+035.000 \\ \hline
2 & P02 & 北边跨 & 330 & 330 & K17+206.000 \\ \hline
3 & P03 & 北边跨 & 501 & 501 & K17+377.000 \\ \hline
4 & P04 & 主跨 & 178 & 838 & K17+714.000 \\ \hline
5 & P05 & 主跨 & 349 & 1009 & K17+885.000 \\ \hline
6 & P06 & 主跨 & 520 & 1180 & K18+056.000 \\ \hline
7 & P07 & 主跨 & 691 & 1351 & K18+227.000 \\ \hline
8 & P08 & 主跨 & 844 & 1504 & K18+380.000 \\ \hline
9 & P09 & 主跨 & 997 & 1657 & K18+533.000 \\ \hline
10 & H1 & 主跨 & 1150 & 1810 & K18+686.000 \\ \hline
11 & H2 & 主跨 & 1303 & 1963 & K18+839.000 \\ \hline
12 & H3 & 主跨 & 1456 & 2116 & K18+992.000 \\ \hline
13 & H4 & 主跨 & 1609 & 2269 & K19+145.000 \\ \hline
14 & H5 & 主跨 & 1780 & 2440 & K19+316.000 \\ \hline
15 & H6 & 主跨 & 1951 & 2611 & K19+487.000 \\ \hline
16 & H7 & 主跨 & 2122 & 2782 & K19+658.000 \\ \hline
17 & H8 & 南侧717 m段 & 178.5 & 3138.5 & K20+014.500 \\ \hline
18 & H9 & 南侧717 m段 & 358.5 & 3318.5 & K20+194.500 \\ \hline
19 & H10 & 南侧717 m段 & 538.5 & 3498.5 & K20+374.500 \\ \hline
20 & H11 & 最南503 m段 & 166 & 3843 & K20+719.000 \\ \hline
21 & H12 & 最南503 m段 & 337 & 4014 & K20+890.000 \\ \hline
\end{longtable}
\endgroup
这些位置由\((159+171+171+159)\)\allowbreak{}、\((178+3\times171+6\times153+3\times171+178)\)\allowbreak{}、\((178.5+180+180+178.5)\)\allowbreak{}和\((166+171+166)\)\allowbreak{}四条跨内尺寸链直接累加，标记为DRW-B。北侧局部倾角可由最终支点/成形线坐标计算，构件质量采用“5769.2 kg核心同型外推”与“10130 kg设计恒载同型外推”两档ASSUMP，直到北侧BOM或称重记录确认；位置、质量和姿态三类信息必须分开标记，不能因后两项未确认而丢弃已知位置。\par
附件2-3较早版本把717 m段尺寸链写成187.5+171+171+187.5 m，对应三道站位为\par
\[
x=3147.5, 3318.5, 3489.5\ \mathrm{m}.
\]
2025施工图与2026复核报告采用178.5+180+180+178.5 m，对应本文主分支的三道站位为\par
\[
x=3138.5, 3318.5, 3498.5\ \mathrm{m}.
\]
这是明确的版本分支，正式计算不得混用；两组位置可作为逐站敏感性，不凭目标频率选择。\par
M05-02给出的H1—H12“横向通道支架与承重索夹角”为：\par
H1 0°；H2 2.82°；H3 5.63°；H4 8.75°；H5 12.44°；H6 15.68°；H7 18.89°；H8 −20.24°；H9 −16.99°；H10 −13.36°；H11 −10.78°；H12 −7.39°。\par
这些角度随主跨上升至南主塔后变号，物理上对应局部纵坡/通道截面姿态。按M05-02的A-A侧向截面，局部通道矩阵应优先绕横桥轴y旋转，而不应绕顺桥轴x旋转；若后续CAD局部轴定义与此不同，以CAD坐标与节点方向余弦为准，绝不依据频率选择旋转轴。对位于局部x-z平面、方向向量为\(\mathbf e=(\cos\beta,\allowbreak 0,\allowbreak \sin\beta)^T\)\allowbreak{}的轴向构件，其秩一刚度包含\par
\[
K_{xz}=k\sin\beta\cos\beta,
\qquad
\frac{K_{xz}}{K_{xx}}=\tan\beta.
\]
当\(|\beta|=20.24^\circ\)\allowbreak{}时，\(|\tan\beta|=0.368721\)\allowbreak{}。这一结果只说明单根秩一轴向构件的局部矩阵存在明显非对角项，尚不能单独证明其对全桥模态显著；全局作用还取决于杆件刚度、端口偏心、安装预应力和对应模态在该站位的相对位移。\par
\phantomsection
\subsection*{2.6. 单幅71道门架的全线站位}
\addcontentsline{toc}{subsection}{2.6. 单幅71道门架的全线站位}
《猫道图纸1225》MD0-01/MD0-02的四跨门架尺寸链可以闭合全部站位。仍以K16+876.000名义四跨北端为\(x=0\)\allowbreak{}，单幅71道门架坐标如下；两幅合计142榀，但两幅各自保留独立自由度。\par
北边跨11道，尺寸链\(45+10\times57+45=660\ \mathrm m\)\allowbreak{}：\par
\[
x=45,102,159,216,273,330,387,444,501,558,615\ \mathrm m.
\]
主跨41道，尺寸链\(64+11\times57+18\times51+11\times57+64=2300\ \mathrm m\)\allowbreak{}：\par
\[
\begin{aligned}
x={}&724,781,838,895,952,1009,1066,1123,1180,1237,1294,1351,\\
&1402,1453,1504,1555,1606,1657,1708,1759,1810,1861,1912,1963,\\
&2014,2065,2116,2167,2218,2269,2326,2383,2440,2497,2554,2611,\\
&2668,2725,2782,2839,2896\ \mathrm m.
\end{aligned}
\]
南侧717 m段11道，尺寸链\(58.5+10\times60+58.5=717\ \mathrm m\)\allowbreak{}：\par
\[
\begin{aligned}
x={}&3018.5,3078.5,3138.5,3198.5,3258.5,3318.5,\\
&3378.5,3438.5,3498.5,3558.5,3618.5\ \mathrm m.
\end{aligned}
\]
最南503 m段8道，图示尺寸链\(52+7\times57+52=503\ \mathrm m\)\allowbreak{}：\par
\[
x=3729,3786,3843,3900,3957,4014,4071,4128\ \mathrm m.
\]
503 m段的说明文字经OCR可疑地出现“52.5 m端距”，但同页标注尺寸链\(5200+7\times5700+5200=50300\)\allowbreak{}与总跨503 m闭合，故站位采用图示52 m端距，并把文字差异登记为版本/OCR冲突。上述71个坐标由实体图纸尺寸链直接推导，标记为DRW-B；派生站位JSON仅用于逐项交叉核对，不能因其中同时含有过时的12道通道清单而整体采信。\par
\phantomsection
\section*{3. 建议采用的五个模型档次}
\addcontentsline{toc}{section}{3. 建议采用的五个模型档次}
\phantomsection
\subsection*{3.1. L0：四跨索系解析骨架}
\addcontentsline{toc}{subsection}{3.1. L0：四跨索系解析骨架}
L0只使用跨长、垂度、重力和边界拓扑，建立四跨浅垂索/悬链线模型，用于检查低阶波速、对称性类别和边跨局部模态量级。它不表示猫道完整扭转，也不用于最终验收。优点是输入最少、推导透明，可以立刻判断某个有限元结果是否出现伪低频、单位错误或错误跨长；缺点是不能包含横梁、面层、门架、横向通道、局部畸变与两幅耦合。\par
\phantomsection
\subsection*{3.2. L1：单幅16索可畸变带状模型}
\addcontentsline{toc}{subsection}{3.2. L1：单幅16索可畸变带状模型}
L1显式保留单幅16根面层承重索，并以离散大/小横梁、面层等效剪切连接和71道门架协调各索。截面不强制刚体，除整体横移v、竖移w和绕顺桥轴扭转\(\theta\)\allowbreak{}外，再保留至少两个畸变坐标：左右面层索带的相对翘曲、门架/扶手系统相对面层的侧摆或转角。L1可以解释单幅横弯、竖弯、扭转和扭转—畸变混合，但不能解释两幅共同/差动分裂与横向通道产生的模态重排。\par
\phantomsection
\subsection*{3.3. L2：双独立猫道子结构＋21道离散横向通道}
\addcontentsline{toc}{subsection}{3.3. L2：双独立猫道子结构＋21道离散横向通道}
L2是推荐的最低可验收模型。上游与下游猫道分别由L1构成，节点、索力、质量、门架、局部坐标和边界全部独立；不施加等位移、等索力、镜像或半模型约束。21道横向通道按站位接入，两端接口至少保留每幅猫道左右两侧的三平动与三转动自由度，通道采用全三维梁/桁架动态刚度或Craig—Bampton降阶。这个档次能够研究共同平移、两幅差动、局部扭转、system-roll、通道弯曲和边跨局部模态，并能分别检验质量-only使有序特征值不升、正半定刚度-only使其不降的严格方向，以及二者同时存在时的竞争与模态重排。\par
\phantomsection
\subsection*{3.4. L3：全三维四跨预应力索—梁—通道有限元模型}
\addcontentsline{toc}{subsection}{3.4. L3：全三维四跨预应力索—梁—通道有限元模型}
L3显式建立两幅共32根面层承重索、扶手索、门架承重索、所有大/小横梁、门架、底梁、面层连接、21道通道以及锚碇/塔顶/辅塔接口；索采用张力激活的几何非线性单元，梁采用三维Timoshenko梁，必要时对开口或薄壁构件加入Vlasov翘曲自由度；局部夹具、鞍座、门架节点和通道端部以转动/滑移弹簧或局部子结构凝聚。先完成自重和安装初应变下的非线性静力平衡，再在线性化切线刚度上求预应力模态。L3用于十四项最终盲验收。\par
\phantomsection
\subsection*{3.5. L4：局部壳/实体子结构回填}
\addcontentsline{toc}{subsection}{3.5. L4：局部壳/实体子结构回填}
若L3的扭转族仍对门架或通道接口高度敏感，则对门架底部、横梁—承重索夹具、通道端部、鞍座和面层网片单元建立局部壳/实体模型，通过静力凝聚或Craig—Bampton获得多端口矩阵，再回填L3。L4不需要全桥使用壳单元，只在扭转能量占比高且接口柔度不确定的部位精细化。\par
\phantomsection
\section*{4. 坐标、自由度与基本符号}
\addcontentsline{toc}{section}{4. 坐标、自由度与基本符号}
采用全桥统一右手坐标：\(x=\text{桩号}-\mathrm{K16{+}876.000}\)\allowbreak{}，顺桥向由北向南为正；\(y\)\allowbreak{}为横桥向，由上游指向下游；\(z\)\allowbreak{}竖直向上。这里\(x=0\)\allowbreak{}是总体布置四跨的名义北端，不是物理锚点。猫道面层承重索的北锚K16+852.105与南锚K21+086.368分别为\(x=-23.895\ \mathrm m\)\allowbreak{}和\(x=4210.368\ \mathrm m\)\allowbreak{}；门架承重索的北锚K16+831.091与南锚K21+101.700则分别为\(x=-44.909\ \mathrm m\)\allowbreak{}和\(x=4225.700\ \mathrm m\)\allowbreak{}；名义四跨端点为\(x=0\)\allowbreak{}和\(x=4180\ \mathrm m\)\allowbreak{}。两幅猫道记为\(p\in\{U,\allowbreak D\}\)\allowbreak{}，每幅纵向钢丝绳记为\(i\)\allowbreak{}，横梁/门架/通道站位记为\(x_j\)\allowbreak{}。第i根索静平衡后的中心线为\(\mathbf r^0_{pi}(s)\)\allowbreak{}，单位切向量为\(\mathbf t^0_{pi}\)\allowbreak{}，静力轴力为\(N^0_{pi}(s)>0\)\allowbreak{}，线质量为\(\mu_{pi}\)\allowbreak{}，表观轴向刚度为\((EA)_{pi}^{\mathrm{app}}\)\allowbreak{}。全局自由度向量记为\(\mathbf q\)\allowbreak{}，静平衡解为\(\mathbf q_0\)\allowbreak{}，模态为\(\boldsymbol\phi_r\)\allowbreak{}，圆频率与频率为\(\omega_r\)\allowbreak{}和\(f_r=\omega_r/(2\pi)\)\allowbreak{}。所有站位JSON、通道坐标和传递矩阵分段均使用该名义原点；若纳入真实锚固/转索段，必须按索族分别增加猫道面层承重索\([-23.895,\allowbreak 0]\)\allowbreak{}、\([4180,\allowbreak 4210.368]\)\allowbreak{}与门架承重索\([-44.909,\allowbreak 0]\)\allowbreak{}、\([4180,\allowbreak 4225.700]\)\allowbreak{}的外伸段，不能继续沿用“from\_north\_anchor”之类易误导字段名，也不能把面层索锚点套给门架索。\par
\phantomsection
\section*{5. L0解析骨架：由图纸垂度得到独立频率尺度}
\addcontentsline{toc}{section}{5. L0解析骨架：由图纸垂度得到独立频率尺度}
\phantomsection
\subsection*{5.1. 等高支点、水平投影均布荷载的抛物线骨架}
\addcontentsline{toc}{subsection}{5.1. 等高支点、水平投影均布荷载的抛物线骨架}
对跨长\(L\)\allowbreak{}、垂度\(h\)\allowbreak{}、按水平投影计的均布重力\(w_x\)\allowbreak{}，水平张力分量\(H\)\allowbreak{}满足平衡方程\(Hz''=w_x\)\allowbreak{}，等高支点解严格为抛物线，并有\par
\[
H=\frac{w_xL^2}{8h}.
\]
若静载质量与横向振动质量均按水平投影计，记为\(\mu_x\)\allowbreak{}，且\(w_x=\mu_xg\)\allowbreak{}，则均匀张弦骨架的波速和两端固定频率为\par
\[
c=\sqrt{\frac{H}{\mu_x}}
=L\sqrt{\frac{g}{8h}},
\qquad
f_n=\frac{n}{2L}c
=n\sqrt{\frac{g}{32h}}.
\]
将猫道成型控制点得到的主跨垂度\(h=227.30\ \mathrm m\)\allowbreak{}和标准重力\(g=9.80665\ \mathrm{m/s^2}\)\allowbreak{}代入，得到不依赖B00、既有MCT及TARGET-FREQ的主跨张力波骨架：\par
\[
f_1=0.036718559\ \mathrm{Hz},\quad
f_2=0.073437118\ \mathrm{Hz},\quad
f_3=0.110155677\ \mathrm{Hz},\quad
f_4=0.146874237\ \mathrm{Hz}.
\]
质量抵消只在静载与惯性采用同一水平投影测度时成立。图纸钢丝绳单位质量通常按实际索长计，必须通过\(\mu_x=\mu_s\,\allowbreak \mathrm ds/\mathrm dx\)\allowbreak{}转换；本主跨同跨同垂度抛物线弧长为2358.570591 m，即\(S/L=1.025465\)\allowbreak{}。若在\(H=w_xL^2/(8h)\)\allowbreak{}中与惯性项同时、同样误用质量测度，质量仍会代数抵消，不能声称必然产生1.25\%的频差；只有当\(H\)\allowbreak{}由独立索力固定、而惯性质量仍把\(\mu_s\)\allowbreak{}误作\(\mu_x\)\allowbreak{}时，才会出现\(\sqrt{S/L}-1\approx1.265\%\)\allowbreak{}的频率尺度偏差。L0只能检查量级，不能把与某一目标频率的剩余差值唯一归因于构件或边界。作为证据回溯对照，误用主缆垂度255.56 m会得到0.034628917 Hz，比正确成型猫道骨架低5.691\%；这一差异足以改变对低阶横弯是否“已经吻合”的判断。\par
\phantomsection
\subsection*{5.2. 名义2300 m L0代理的精确自重悬链线}
\addcontentsline{toc}{subsection}{5.2. 名义2300 m L0代理的精确自重悬链线}
原稿此处把“沿水平投影均布荷载”与双曲余弦悬链线混写。沿水平投影均布的恒定荷载仍满足\(Hz''=w_x\)\allowbreak{}并产生抛物线；双曲余弦悬链线对应沿当前索弧长均布的自重\(w_s\)\allowbreak{}，其平衡方程和参数为\par
\[
Hz''=w_s\sqrt{1+z'^2},
\qquad
a=\frac{H}{w_s},
\qquad
z(x)=a\cosh\left(\frac{x-x_c}{a}\right)+c_0.
\]
两端高程条件为\par
\[
z(0)=z_A,\qquad z(L)=z_B,
\]
弧长为\par
\[
S=a\left[
\sinh\left(\frac{L-x_c}{a}\right)
+
\sinh\left(\frac{x_c}{a}\right)
\right].
\]
当两端等高时\(x_c=L/2\)\allowbreak{}，垂度满足\par
\[
h=a\left[
\cosh\left(\frac{L}{2a}\right)-1
\right].
\]
为与5.1节的名义2300 m骨架保持同一代理口径并单独观察抛物线—悬链线差异，本节有意把名义跨长\(L=2300\ \mathrm m\)\allowbreak{}与成型控制高差\(h=227.30\ \mathrm m\)\allowbreak{}组合成L0代理，得\(a=2946.273677\ \mathrm m\)\allowbreak{}、\(S=2358.848281\ \mathrm m\)\allowbreak{}。这不是实际鞍点间几何；21.4节另以实际鞍点间\(L_{\mathrm{eff}}=2286.642\ \mathrm m\)\allowbreak{}拟合成型曲率，两个分支不交叉使用坐标或形函数。在固定端面外振动且质量沿索长均布时，本L0代理的振型满足\par
\[
-H\psi''(x)
=\omega^2\mu_s
\cosh\left(\frac{x-L/2}{a}\right)\psi(x),
\qquad H=\mu_sga.
\]
独立的64项正弦Ritz解给出前四阶\(0.036768481\)\allowbreak{}、\(0.073114885\)\allowbreak{}、\(0.109560267\)\allowbreak{}、\(0.146028469\ \mathrm{Hz}\)\allowbreak{}，相对抛物线均匀张弦骨架分别变化\(+0.1360\%\)\allowbreak{}、\(-0.4388\%\)\allowbreak{}、\(-0.5405\%\)\allowbreak{}、\(-0.5758\%\)\allowbreak{}。320单元变质量有限元与Ritz解前四阶最大差小于0.01\%。该修正并非统一缩放；面内竖弯仍需要实际\(EA_t\)\allowbreak{}、无应力长度和静力平衡状态。正式四跨模型应使用真实支点坐标、索的无应力长度或安装初应变、离散自重和支点接触条件求解，而不把所有跨强制为等高抛物线或悬链线。\par
\phantomsection
\subsection*{5.3. 面内竖弯的浅索兼容模型}
\addcontentsline{toc}{subsection}{5.3. 面内竖弯的浅索兼容模型}
现有材料已经给出一条不读取MCT索力即可执行的面内验证路径。复核报告表1-1给单幅均布恒载\(q_x=2.766\ \mathrm{kN/m}\)\allowbreak{}，第6页给十五根普通\(\phi50\)\allowbreak{}绳和一根智慧芯\(\phi50\)\allowbreak{}绳的设计\(E\)\allowbreak{}与金属面积，得到\par
\[
(EA)_{16}
=120\ \mathrm{GPa}\times
\left(15\times1400.42+1292.45\right)\ \mathrm{mm^2}
=2.675850\ \mathrm{GN}.
\]
在两端轴向锁定、忽略轴向惯性、集中质量和四跨鞍座作用的浅索近似下，面内竖向位移\(w(x,\allowbreak t)\)\allowbreak{}的二阶能量为\par
\[
U_w=\frac12H\int_0^L w'^2\,\mathrm dx
+\frac12\frac{(EA)_{16}}{S}
\left(\int_0^L z_0'w'\,\mathrm dx\right)^2,
\qquad
T_w=\frac12\mu_x\int_0^L\dot w^2\,\mathrm dx.
\]
第一项是水平分力几何刚度，第二项是全跨伸长兼容产生的秩一刚度。为与5.1—5.2节数值对照保持一致，这里仍明确使用名义2300 m L0代理：将成型控制高差227.300 m与名义跨长组合，得\(H=q_xL^2/(8h)=8.046711\ \mathrm{MN}\)\allowbreak{}，并取\(\mu_x=q_x/g=282.0535\ \mathrm{kg/m}\)\allowbreak{}。该数值不是实际2286.642 m鞍间静力解；统一实际几何只在待求解的L1/L2分支中使用。采用64项正弦Galerkin基并在完成预测后打开四个V族目标，得到：\par
VA1：0.073437118 Hz，相对0.0700 Hz为+4.910\%；\par
VS1：0.101932295 Hz，相对0.1028 Hz为−0.844\%；\par
VA2：0.146874237 Hz，相对0.1438 Hz为+2.138\%；\par
VS2：0.159158688 Hz，相对0.1744 Hz为−8.739\%。\par
四项平均绝对相对误差为4.158\%。反对称竖弯不激发全跨一阶长度变化，因此VA1、VA2分别落在第2、4张弦谐波；对称竖弯由秩一兼容刚度重排。VS1在未调参条件下接近目标，与轴向兼容机制和设计\(EA\)\allowbreak{}量级相容，但单一频率不能独立识别或证明二者；VS2偏低则明确指向当前简化未包含的集中质量/刚度、下拉系统、四跨边界和工作点切线\(EA_t\)\allowbreak{}，不能用单一比例修正。报告中的索束最大轴力只作为冻结后对照，不进入上述\(H\)\allowbreak{}计算。\par
\phantomsection
\section*{6. 三维索的非线性静力、切线刚度与预应力模态}
\addcontentsline{toc}{section}{6. 三维索的非线性静力、切线刚度与预应力模态}
\phantomsection
\subsection*{6.1. 总势能与静平衡}
\addcontentsline{toc}{subsection}{6.1. 总势能与静平衡}
对第i根索，以一般参考弧长坐标\(S\in[0,\allowbreak L_R]\)\allowbreak{}描述当前中心线\(\mathbf r(S)\)\allowbreak{}，伸长比与相对该参考状态的工程应变定义为\par
\[
\lambda=\left\|\frac{\partial\mathbf r}{\partial S}\right\|,
\qquad
\varepsilon=\lambda-1-\varepsilon_0,
\]
其中\(\varepsilon_0\)\allowbreak{}仅表示相对一般参考状态的制造/张拉初应变。若\(S\)\allowbreak{}已取无应力参考，则必须令\(\varepsilon_0=0\)\allowbreak{}并回到\(\varepsilon=\lambda-1\)\allowbreak{}，不得重复扣除初应变。若钢丝绳的表观刚度随荷载区间、预循环和加载路径变化，应以单位参考长度储能\(W(\lambda;\allowbreak \varepsilon_0)\)\allowbreak{}或经试验识别的增量本构表示：\par
\[
U_c=\int_0^{L_R}W(\lambda;\varepsilon_0)\,\mathrm dS,
\qquad
N=\frac{\partial W}{\partial\lambda},
\qquad
(EA)_t=\frac{\partial N}{\partial\lambda}.
\]
只有在指定工作区间内可近似线弹性时，才可把\(W\)\allowbreak{}写成二次型；模态阶段使用静平衡点的卸载/再加载切线\((EA)_t\)\allowbreak{}，不得把一个工作点切线常数无条件外推到完整张拉历程。\par
全体系总势能写成\par
\[
\Pi(\mathbf q)=
\sum_{\mathrm{cables}}U_c
+
\sum_{\mathrm{beams}}U_b
+
\sum_{\mathrm{connections}}U_k
-
\mathbf q^T\mathbf f_g
-
\mathbf q^T\mathbf f_{\mathrm{other}}.
\]
静平衡满足\par
\[
\delta\Pi(\mathbf q_0)=\mathbf 0,
\qquad
\mathbf R(\mathbf q_0)=\mathbf f_{\mathrm{int}}(\mathbf q_0)-\mathbf f_{\mathrm{ext}}=\mathbf0.
\]
必须先求\(\mathbf q_0\)\allowbreak{}，随后才能构造预应力切线刚度；直接在未平衡图纸几何上求模态会把重力不平衡和错误初应力带入低阶频率。\par
\phantomsection
\subsection*{6.2. 连续体二阶变分}
\addcontentsline{toc}{subsection}{6.2. 连续体二阶变分}
在静平衡构形附近，索的位移扰动为\(\mathbf u\)\allowbreak{}，单位切向量为\(\mathbf t^0\)\allowbreak{}，横向投影算子为\par
\[
\mathbf P=\mathbf I-\mathbf t^0\mathbf t^{0T}.
\]
以参考弧长\(S\)\allowbreak{}为积分坐标，二阶势能可写成\par
\[
\Pi_i^{(2)}=
\frac12\int_0^{L_R}
\left[
(EA)_{t,i}\left(\mathbf t_i^0\cdot\mathbf u_{i,S}\right)^2
+
\frac{N_i^0}{\lambda_i}
\mathbf u_{i,S}^{T}\mathbf P_i\mathbf u_{i,S}
\right]\mathrm dS.
\]
若改用当前弧长\(s\)\allowbreak{}，材料项系数相应为\((EA)_t\lambda\)\allowbreak{}，几何项为\(N\)\allowbreak{}。第一项是材料轴向刚度，第二项是静张力产生的几何刚度。扭转与面外横弯的低频骨架主要来自第二项；面内竖弯还包含静态曲率引起的轴向—法向耦合，不能直接视为与面外横弯完全相同。原稿的当前弧长简式只在\(\lambda\approx1\)\allowbreak{}的小弹性应变近似下成立。\par
\phantomsection
\subsection*{6.3. 两节点三维索/桁单元切线矩阵}
\addcontentsline{toc}{subsection}{6.3. 两节点三维索/桁单元切线矩阵}
对当前长度l、无应力长度\(L_0\)\allowbreak{}、方向向量\(\mathbf n\)\allowbreak{}、轴力N的两节点单元，定义\par
\[
k_a=\frac{(EA)_t}{L_0},
\qquad
\mathbf A_e=k_a\mathbf n\mathbf n^T
+
\frac{N}{l}\left(\mathbf I-\mathbf n\mathbf n^T\right).
\]
其六自由度切线刚度为\par
\[
\mathbf K_e^T=
\begin{bmatrix}
\mathbf A_e&-\mathbf A_e\\
-\mathbf A_e&\mathbf A_e
\end{bmatrix}.
\]
若索线质量按参考长度为\(\mu_0\)\allowbreak{}，一致质量矩阵为\par
\[
\mathbf M_e=
\frac{\mu_0L_0}{6}
\begin{bmatrix}
2\mathbf I&\mathbf I\\
\mathbf I&2\mathbf I
\end{bmatrix}.
\]
索单元采用仅受拉激活条件；若迭代中N≤0，应进入松索活动集或使用平滑张力截断，不能保留负轴力形成虚假压杆刚度。\par
\phantomsection
\section*{7. 单幅多索截面：整体平移、扭转与畸变的统一推导}
\addcontentsline{toc}{section}{7. 单幅多索截面：整体平移、扭转与畸变的统一推导}
\phantomsection
\subsection*{7.1. 刚体截面运动}
\addcontentsline{toc}{subsection}{7.1. 刚体截面运动}
先取单幅猫道某一纵向截面。第i根索相对参考点的坐标为\((y_i,\allowbreak z_i)\)\allowbreak{}，截面整体横移、竖移和绕x轴转角分别为\(v(x),\allowbreak w(x),\allowbreak \theta(x)\)\allowbreak{}。小转角下第i根索的横向和竖向位移为\par
\[
u_{y,i}=v-z_i\theta,
\qquad
u_{z,i}=w+y_i\theta.
\]
为了允许左右索带、扶手系统和门架产生相对畸变，扩展为\par
\[
u_{y,i}=v-z_i\theta+
\sum_{\alpha=1}^{n_d}\psi^{v}_{i\alpha}\chi_\alpha,
\qquad
u_{z,i}=w+y_i\theta+
\sum_{\alpha=1}^{n_d}\psi^{w}_{i\alpha}\chi_\alpha,
\]
其中\(\chi_\alpha(x)\)\allowbreak{}为畸变坐标，\(\psi_{i\alpha}\)\allowbreak{}可由一榀横梁/门架局部模型的截面特征模态得到，也可先选取具有明确物理含义的基：左右8索带相对翘曲、扶手/门架相对面层侧摆、门架开口剪切。若只保留\(v,\allowbreak w,\allowbreak \theta\)\allowbreak{}，截面会被强制为刚体，可能高估面层协调刚度；若完全逐索自由而不建横梁/网片，又会低估扭转和产生非物理独立索模态。\par
\phantomsection
\subsection*{7.2. 张力几何刚度展开与不对称耦合}
\addcontentsline{toc}{subsection}{7.2. 张力几何刚度展开与不对称耦合}
在浅垂度、小转角并以水平投影坐标\(x\)\allowbreak{}积分的约化中，第i根索应使用工作态水平张力分量\(H_i^0\)\allowbreak{}，其横向几何能为\par
\[
U_{g,i}=\frac12\int H_i^0
\left[
\left(v'-z_i\theta'\right)^2
+
\left(w'+y_i\theta'\right)^2
\right]\mathrm dx.
\]
求和并定义\par
\[
H_\Sigma=\sum_iH_i^0,
\quad
S_y^{H}=\sum_iH_i^0y_i,
\quad
S_z^{H}=\sum_iH_i^0z_i,
\quad
J_H=\sum_iH_i^0(y_i^2+z_i^2),
\]
得到\par
\[
U_g=\frac12\int
\left[
H_\Sigma(v'^2+w'^2)
-2S_z^{H}v'\theta'
+2S_y^{H}w'\theta'
+J_H\theta'^2
\right]\mathrm dx.
\]
完全对称且成对水平分力相同的理想截面可能有\(S_y^{H}=S_z^{H}=0\)\allowbreak{}，横移、竖移与扭转在这一阶近似中解耦。实际模型中，左右索坐标误差、安装索力差、扶手/门架承重索的高度偏心、门架类型差异、施工附加质量和通道端部力会使\(S_y^{H}\)\allowbreak{}或\(S_z^{H}\)\allowbreak{}非零，于是平移—扭转耦合自然出现。正确处理方式是逐索保存\(y_i,\allowbreak z_i,\allowbreak H_i\)\allowbreak{}，不人为设置耦合系数，也不通过频率拟合耦合一阶矩。对大垂度、不等高或竖向面内运动，正式模型必须回到工作态弧长坐标\(s\)\allowbreak{}的二阶变分，保留切向—法向耦合；上式只作为浅索横向骨架，不能把随弧长变化的全轴力\(N_i(s)\)\allowbreak{}直接代入\(x\)\allowbreak{}积分。\par
\phantomsection
\subsection*{7.3. 截面质量矩阵与偏心质量}
\addcontentsline{toc}{subsection}{7.3. 截面质量矩阵与偏心质量}
在以水平投影坐标\(x\)\allowbreak{}积分的约化模型中，第a个按水平投影计的单位长度质量\(\mu_{x,\allowbreak a}\)\allowbreak{}位于\((y_a,\allowbreak z_a)\)\allowbreak{}，其动能为\par
\[
T_a=\frac12\mu_{x,a}
\left[
(\dot v-z_a\dot\theta)^2
+(\dot w+y_a\dot\theta)^2
\right].
\]
定义\par
\[
\mu_\Sigma=\sum_a\mu_{x,a},
\quad
S_{my}=\sum_a\mu_{x,a}y_a,
\quad
S_{mz}=\sum_a\mu_{x,a}z_a,
\quad
I_{m\theta}=\sum_a\mu_{x,a}(y_a^2+z_a^2)+\sum_aJ_{p,x,a},
\]
则\([v,\allowbreak w,\allowbreak \theta]^T\)\allowbreak{}对应的单位长度质量矩阵为\par
\[
\mathbf M_s=
\begin{bmatrix}
\mu_\Sigma&0&-S_{mz}\\
0&\mu_\Sigma&S_{my}\\
-S_{mz}&S_{my}&I_{m\theta}
\end{bmatrix}.
\]
门架高7.2—9.6 m，门架承重索与上扶手索位于面层上方，哪怕其质量占比不大，\(z^2\)\allowbreak{}放大后也可能显著影响\(I_{m\theta}\)\allowbreak{}和\(S_{mz}\)\allowbreak{}。因此把所有附属质量投到面层中心节点会破坏扭转惯量和横移—扭转耦合。\par
\phantomsection
\subsection*{7.4. 张力索束的“平移—扭转退化”及其意义}
\addcontentsline{toc}{subsection}{7.4. 张力索束的“平移—扭转退化”及其意义}
若所有纵向索满足相同波速比\par
\[
\frac{H_i^0}{\mu_{x,i}}=c^2,
\]
且暂不计横梁、门架、面层和通道的附加质量/刚度，并令各索自身关于其质心的极惯量项\(J_{p,\allowbreak x,\allowbreak i}=0\)\allowbreak{}，定义纯平移质量形成的扭转二次矩\par
\[
I_{m\theta}^{(0)}=\sum_i\mu_{x,i}r_i^2,
\qquad r_i^2=y_i^2+z_i^2,
\]
则\par
\[
J_H=\sum_iH_i^0r_i^2
=c^2\sum_i\mu_{x,i}r_i^2
=c^2I_{m\theta}^{(0)}.
\]
于是纯张力索束的扭转波速\par
\[
c_\theta=\sqrt{\frac{J_H}{I_{m\theta}^{(0)}}}=c
\]
与平移弦波速度相同。这里\(\mu_{x,\allowbreak i}\)\allowbreak{}必须与\(H_i\)\allowbreak{}同为水平投影口径；若图纸质量按实际索长给出，则先用\(\mu_{x,\allowbreak i}=\mu_{s,\allowbreak i}\,\allowbreak \mathrm ds/\mathrm dx\)\allowbreak{}转换。这个结果说明：只把具有同一\(H_i/\mu_{x,\allowbreak i}\)\allowbreak{}的多根索刚性捆在一起，会让扭转和平移共享同一频率骨架；边界条件也必须对两类场一致，截面自身极惯量\(J_{p,\allowbreak x,\allowbreak i}\)\allowbreak{}等额外项也须暂不计入。真实分支分裂可来自离散门架与横向通道的相对转动刚度、面层/横梁畸变、索间波速差、四跨边界和偏心集中质量。独立数值检验将16根同波速索组装后，前四阶平移与扭转频率在机器精度内重合，验证了该退化判据；这只是合成单元测试，不是项目索力识别结果。\par
\phantomsection
\subsection*{7.5. 连续扭转方程与Rayleigh商}
\addcontentsline{toc}{subsection}{7.5. 连续扭转方程与Rayleigh商}
将单幅猫道约化为扭转场\(\theta(x,\allowbreak t)\)\allowbreak{}，令\(C_\theta(x)\)\allowbreak{}包含索张力二次矩和连续构件的Saint-Venant贡献。内部面层、门架和横向通道必须对刚体运动保持客观性，因此只惩罚相连端口之间的相对位移/转角；只有真实连接到地基或塔体的支承才可写成绝对转角弹簧。若两幅端口采用完全相同的物理观测坐标，标量相对转角才可退化写成\par
\[
U_{j}^{\mathrm{rel}}
=\frac12k_{\theta,j}
\left[\theta_U(x_j)-\theta_D(x_j)\right]^2.
\]
真实通道的完整端口关系必须写成\par
\[
\mathbf d_{j,\mathrm{rel}}
=\mathbf B_{U,j}\mathbf q_U-\mathbf B_{D,j}\mathbf q_D,
\]
\[
\mathbf K_{p,j}=
\begin{bmatrix}
\mathbf B_{U,j}^{T}\mathbf K_j\mathbf B_{U,j}&
-\mathbf B_{U,j}^{T}\mathbf K_j\mathbf B_{D,j}\\
-\mathbf B_{D,j}^{T}\mathbf K_j\mathbf B_{U,j}&
\mathbf B_{D,j}^{T}\mathbf K_j\mathbf B_{D,j}
\end{bmatrix},
\]
并有\(U_j=\tfrac12\mathbf d_{j,\allowbreak \mathrm{rel}}^T\mathbf K_j\mathbf d_{j,\allowbreak \mathrm{rel}}\)\allowbreak{}。两幅猫道内侧连接点的横向偏心符号相反，例如\(u_z^U=w_U+a\theta_U\)\allowbreak{}、\(u_z^D=w_D-a\theta_D\)\allowbreak{}，故相对竖移为\((w_U-w_D)+a(\theta_U+\theta_D)\)\allowbreak{}：同号局部扭转也可能使通道受力，反号局部扭转反而可能释放。只有\(\mathbf B_{U,\allowbreak j}=\mathbf B_{D,\allowbreak j}\)\allowbreak{}时，端口刚度才退化为纯差动块。所有\(\mathbf B\)\allowbreak{}映射都须通过六个全体系刚体平移/转动的零应变检验。连续单幅无地基约束时，主算子为\par
\[
I_{m\theta}(x)\ddot\theta
-\frac{\partial}{\partial x}
\left[C_\theta(x)\theta'\right]
=\text{平移、畸变与端口相对作用项}.
\]
若试函数包含两幅猫道，Rayleigh商的分子应按\(\sum_j\mathbf d_{j,\allowbreak \mathrm{rel}}^T\mathbf K_j\mathbf d_{j,\allowbreak \mathrm{rel}}\)\allowbreak{}组装，分母加入由同一端口映射得到的通道一致质量。Vlasov翘曲对Rayleigh分子的贡献应另写成\(\int EC_w(\theta'')^2\,\allowbreak \mathrm dx\)\allowbreak{}；这里是刚度二次型分子，因此不再写能量式中的\(1/2\)\allowbreak{}，且不能混入一阶导数系数\(C_\theta\)\allowbreak{}。通道究竟主要作用于共同还是差动坐标，必须在真实\(\mathbf B_U,\allowbreak \mathbf B_D\)\allowbreak{}下变换后判断，不能仅按两幅中心线转角的同号/反号预判。必须分别进行“仅加端口一致质量”“仅加客观刚度”“质量和刚度同时加入”三组消融，并用模态能量解释结果。\par
\phantomsection
\section*{8. 横梁、面层、门架和局部接口的建模}
\addcontentsline{toc}{section}{8. 横梁、面层、门架和局部接口的建模}
\phantomsection
\subsection*{8.1. Timoshenko梁能量}
\addcontentsline{toc}{subsection}{8.1. Timoshenko梁能量}
大横梁、小横梁、门架梁柱和通道杆件采用三维Timoshenko梁时，其一维能量可写为\par
\[
U_b=\frac12\int
\left[
EA\varepsilon_x^2
+k_{sy}GA\gamma_y^2
+k_{sz}GA\gamma_z^2
+EI_y\chi_y^2
+EI_z\chi_z^2
+GJ_t\vartheta'^2
\right]\mathrm dx.
\]
对薄壁开口或连接导致明显翘曲的构件，可加入Vlasov项\par
\[
U_{\mathrm{warp}}=\frac12\int EC_w\vartheta''^2\,\mathrm dx.
\]
方矩管初步截面面积和惯性矩可由\par
\[
A=BH-(B-2t)(H-2t),
\]
\[
I_y=\frac{BH^3-(B-2t)(H-2t)^3}{12},
\]
闭口薄壁管的Saint-Venant扭转常数可用Bredt近似\par
\[
J_t\approx\frac{4A_m^2}{\displaystyle\oint\frac{\mathrm ds}{t}}.
\]
正式截面性质应优先采用型钢产品标准、制造图或截面库中考虑圆角的数值；质量以图纸BOM为控制，避免几何理想化与表中单重不一致。\par
\phantomsection
\subsection*{8.2. 面层钢丝网与索间连接}
\addcontentsline{toc}{subsection}{8.2. 面层钢丝网与索间连接}
面层钢丝网对低阶整体模态主要提供索间剪切协调和分布质量，不宜直接设成无限刚性膜。可采用正交膜/网格模型，或在相邻承重索之间设置经独立拉剪试验确定的等效连接。任一连接的相对变形必须由两连接点在实际连接方向上的伸长、剪切和弯曲得到，并通过任意刚体平移和刚体转动下\(U=0\)\allowbreak{}的客观性检验。以沿工作态单位方向\(\mathbf e_{ik}\)\allowbreak{}的轴向连接为例，离散能量为\par
\[
U_{ik}^{\mathrm{ax}}
=\frac12\int k_{ik}^{\mathrm{ax}}
\left[\mathbf e_{ik}^{T}
(\mathbf u_i-\mathbf u_k)\right]^2\,\mathrm dx.
\]
剪切和弯曲通道按同一局部基构造，不能用三个全局笛卡尔方向的位移差同时充当互不相关弹簧，因为这种写法会对刚体转动产生虚假能量。\(k_{ik}^{\mathrm{ax}}\)\allowbreak{}及局部剪切/弯曲刚度不能由十四项频率拟合，应由网片规格、连接间距及一段实物面层的拉剪/翘曲试验确定。首轮敏感性可设置“无面层剪切、试验下限、试验均值、试验上限”四档。\par
\phantomsection
\subsection*{8.3. 大/小横梁的离散站位}
\addcontentsline{toc}{subsection}{8.3. 大/小横梁的离散站位}
M00-04说明主跨相邻大横梁约9 m、其间两道小横梁形成3 m节距；边跨大横梁约6 m、其间一道小横梁形成3 m节距；踏步方木约0.5 m。正式模型按站位离散，不把横梁连续涂抹。横梁与16根承重索的连接可分三级：铰接平动协调、带转动弹簧的半刚接、局部夹具子结构。每级均保持质量不变，只改变独立证据支持的连接刚度。\par
\phantomsection
\subsection*{8.4. 门架必须保留转动端口}
\addcontentsline{toc}{subsection}{8.4. 门架必须保留转动端口}
门架高度7.2—9.6 m，顶部连接门架承重索和施工设备，底部连接面层/底梁，具有显著的横向弯曲、开口剪切和局部转动。建议每榀门架建立三维梁模型，接口至少包括底部左、底部右、顶部左、顶部右四个端口，每端口六自由度。若只保留端口平动，横剪和弯曲的客观转动关系会丢失，凝聚矩阵可能只剩轴向或出现虚假刚度。门架静力凝聚写成\par
\[
\begin{bmatrix}
\mathbf K_{bb}&\mathbf K_{bi}\\
\mathbf K_{ib}&\mathbf K_{ii}
\end{bmatrix}
\begin{bmatrix}
\mathbf u_b\\
\mathbf u_i
\end{bmatrix}
=
\begin{bmatrix}
\mathbf f_b\\
\mathbf0
\end{bmatrix},
\]
消去内部自由度得到\par
\[
\bar{\mathbf K}
=\mathbf K_{bb}
-\mathbf K_{bi}\mathbf K_{ii}^{-1}\mathbf K_{ib}.
\]
若需要保留门架自身低阶动力，采用下一节的Craig—Bampton矩阵，不能只做静力凝聚。\par
\phantomsection
\section*{9. 横向通道的三级等效方式}
\addcontentsline{toc}{section}{9. 横向通道的三级等效方式}
\phantomsection
\subsection*{9.1. P-A：秩一轴向连接，仅用于极低阶筛查}
\addcontentsline{toc}{subsection}{9.1. P-A：秩一轴向连接，仅用于极低阶筛查}
若只研究某一根连接杆的轴向作用，端点相对位移\(\Delta\mathbf u\)\allowbreak{}沿单位方向\(\mathbf e\)\allowbreak{}的伸长为\(\delta=\mathbf e^T\Delta\mathbf u\)\allowbreak{}，能量为\par
\[
U=\frac12k\delta^2
=\frac12\Delta\mathbf u^T
\left(k\mathbf e\mathbf e^T\right)
\Delta\mathbf u.
\]
两端组装矩阵为\par
\[
\mathbf K_{\mathrm{link}}=k
\begin{bmatrix}
\mathbf e\mathbf e^T&-\mathbf e\mathbf e^T\\
-\mathbf e\mathbf e^T&\mathbf e\mathbf e^T
\end{bmatrix}.
\]
这一档只能检查连接拓扑与方向，无法代表49.655 m三角桁架的弯曲、扭转、端部偏心和自身动力，不能用于扭转族验收。\par
\phantomsection
\subsection*{9.2. P-B：上下弦/多端口等效梁格}
\addcontentsline{toc}{subsection}{9.2. P-B：上下弦/多端口等效梁格}
把通道端部在每幅猫道上至少分成上、下两个接口，保留端口平动和转动；上下弦、斜腹杆和横联形成力偶，可抵抗两幅猫道差动竖移和局部扭转。端口运动\(\mathbf d_j\)\allowbreak{}的能量写成\par
\[
U_j=\frac12\mathbf d_j^T\mathbf K_j\mathbf d_j,
\qquad
T_j=\frac12\dot{\mathbf d}_j^T\mathbf M_j\dot{\mathbf d}_j.
\]
此档可以用图纸杆件截面直接组装，无需目标频率校准，适合快速比较通道质量与刚度对TA/TS的影响。\par
\phantomsection
\subsection*{9.3. P-C：全通道Craig—Bampton降阶，推荐用于最终模型}
\addcontentsline{toc}{subsection}{9.3. P-C：全通道Craig—Bampton降阶，推荐用于最终模型}
将通道自由度分为接口b和内部i。固定接口内部模态满足\par
\[
\mathbf K_{ii}\boldsymbol\Phi
=
\mathbf M_{ii}\boldsymbol\Phi\boldsymbol\Omega^2.
\]
约束模态为\par
\[
\boldsymbol\Psi_c=-\mathbf K_{ii}^{-1}\mathbf K_{ib}.
\]
采用变换\par
\[
\begin{bmatrix}
\mathbf u_b\\
\mathbf u_i
\end{bmatrix}
=
\underbrace{
\begin{bmatrix}
\mathbf I&\mathbf0\\
\boldsymbol\Psi_c&\boldsymbol\Phi
\end{bmatrix}}_{\mathbf T_{CB}}
\begin{bmatrix}
\mathbf u_b\\
\boldsymbol\eta
\end{bmatrix},
\]
得到\par
\[
\mathbf K_r=\mathbf T_{CB}^T\mathbf K\mathbf T_{CB},
\qquad
\mathbf M_r=\mathbf T_{CB}^T\mathbf M\mathbf T_{CB}.
\]
接口至少保留两幅猫道端部所有真实连接点的三平动与三转动。冻结前不读取目标频率；固定接口模态截断先依据L0独立骨架和计划求解带宽设为1 Hz，再以0.5、1.0、2.0 Hz三档检查前60阶全局特征对的收敛。冻结并打开TARGET-FREQ后，只把目标带宽用于验收截断充分性，不能回调局部参数。通道质量必须只计一次，护栏、网片和端部附件按BOM加入同一局部模型。\par
\phantomsection
\subsection*{9.4. 逐站坐标变换}
\addcontentsline{toc}{subsection}{9.4. 逐站坐标变换}
第j道通道局部坐标到全局坐标的主动旋转矩阵记为\(\mathbf R_j\)\allowbreak{}，其第一列必须等于该通道局部\(x\)\allowbreak{}轴在全局坐标中的方向。按当前图纸A-A定义，若\(\mathbf e_j=(\cos\beta_j,\allowbreak 0,\allowbreak \sin\beta_j)^T\)\allowbreak{}，可写成\par
\[
\mathbf R_y(\beta_j)=
\begin{bmatrix}
\cos\beta_j&0&-\sin\beta_j\\
0&1&0\\
\sin\beta_j&0&\cos\beta_j
\end{bmatrix}.
\]
对所有端口平动和转动构造局部到全局的块对角映射\(\mathbf T_j\)\allowbreak{}，令\(\mathbf d^G=\mathbf T_j\mathbf d^L\)\allowbreak{}，则\par
\[
\mathbf K_j^{G}=\mathbf T_j\mathbf K_j^{L}\mathbf T_j^T,
\qquad
\mathbf M_j^{G}=\mathbf T_j\mathbf M_j^{L}\mathbf T_j^T.
\]
H1—H12逐站采用图纸角度；北侧站位采用北侧详图或实测姿态。将所有通道统一成水平同一矩阵会丢失塔区变号和纵向—竖向耦合。\par
\phantomsection
\section*{10. 双幅猫道、共同/差动坐标与不完美对称}
\addcontentsline{toc}{section}{10. 双幅猫道、共同/差动坐标与不完美对称}
\phantomsection
\subsection*{10.1. 双子结构方程}
\addcontentsline{toc}{subsection}{10.1. 双子结构方程}
上、下游猫道的独立刚度与质量记为\(\mathbf K_U,\allowbreak \mathbf K_D,\allowbreak \mathbf M_U,\allowbreak \mathbf M_D\)\allowbreak{}。第j道通道的端口映射按7.5节写成\(\mathbf B_{U,\allowbreak j},\allowbreak \mathbf B_{D,\allowbreak j}\)\allowbreak{}，其全局刚度块为\par
\[
\mathbf K_{p,j}=
\begin{bmatrix}
\mathbf B_{U,j}^{T}\mathbf K_j\mathbf B_{U,j}&
-\mathbf B_{U,j}^{T}\mathbf K_j\mathbf B_{D,j}\\
-\mathbf B_{D,j}^{T}\mathbf K_j\mathbf B_{U,j}&
\mathbf B_{D,j}^{T}\mathbf K_j\mathbf B_{D,j}
\end{bmatrix}.
\]
完整体系写成\par
\[
\mathbf K=
\begin{bmatrix}
\mathbf K_U&\mathbf0\\
\mathbf0&\mathbf K_D
\end{bmatrix}
+\sum_j\mathbf K_{p,j},
\qquad
\mathbf M=
\begin{bmatrix}
\mathbf M_U&\mathbf0\\
\mathbf0&\mathbf M_D
\end{bmatrix}
+\sum_j\mathbf M_{p,j},
\]
其中\(\mathbf M_{p,\allowbreak j}\)\allowbreak{}由同一真实端口运动映射和通道内部自由度凝聚得到。只有在两幅端口同构且\(\mathbf B_U=\mathbf B_D\)\allowbreak{}时，才可把通道简写成单一\(\mathbf K_c\)\allowbreak{}的\([+,\allowbreak -;\allowbreak -,\allowbreak +]\)\allowbreak{}块。\par
定义共同与差动坐标\par
\[
\mathbf q_+=\frac{\mathbf q_U+\mathbf q_D}{\sqrt2},
\qquad
\mathbf q_-=\frac{\mathbf q_U-\mathbf q_D}{\sqrt2},
\]
对两幅自身矩阵定义\par
\[
\bar{\mathbf K}=\frac{\mathbf K_U+\mathbf K_D}{2},
\qquad
\Delta\mathbf K=\frac{\mathbf K_U-\mathbf K_D}{2}.
\]
若暂取上述同构退化且通道不含未配对内部自由度，刚度变换后才有\par
\[
\mathbf K_{\pm}=
\begin{bmatrix}
\bar{\mathbf K}&\Delta\mathbf K\\
\Delta\mathbf K&\bar{\mathbf K}+2\mathbf K_c
\end{bmatrix}.
\]
质量矩阵即使不含通道也具有\par
\[
\widetilde{\mathbf M}_{\pm}=
\begin{bmatrix}
\bar{\mathbf M}&\Delta\mathbf M\\
\Delta\mathbf M&\bar{\mathbf M}
\end{bmatrix}
+\widetilde{\mathbf M}_{p},
\qquad
\bar{\mathbf M}=\frac{\mathbf M_U+\mathbf M_D}{2},
\quad
\Delta\mathbf M=\frac{\mathbf M_U-\mathbf M_D}{2}.
\]
因此真实动能含交叉项\(2\dot{\mathbf q}_{+}^{T}\Delta\mathbf M\dot{\mathbf q}_{-}\)\allowbreak{}，再加通道内部质量造成的未配对项；不能令共同与差动参与率自动互补。当两幅完全相同、端口映射具有正确镜像关系且通道理想对称时，两类模态才可能解耦。实际几何、索力、质量、门架、通道端部和施工状态使\(\Delta\mathbf K\)\allowbreak{}、\(\Delta\mathbf M\)\allowbreak{}或端口映射差异非零，两类模态发生混合与避交。因此模型不应先按对称/反对称拆开求解，正确做法是求全体系，再用带符号交换算子和真实端口能量解释振型。\par
\phantomsection
\subsection*{10.2. 三类不对称必须分开记录}
\addcontentsline{toc}{subsection}{10.2. 三类不对称必须分开记录}
纵向不对称来自660—2300—717—503 m跨序和最南辅助塔；横向不对称来自两幅实测坐标、索力和质量差，以及每幅16根承重索中1根智慧芯绳的具体横向位置；离散构造不对称来自H1—H12倾角、M1/M2门架、非均匀站距和北/南通道细节差异。首轮若尚无两幅实测差异，可让两幅名义构件参数相同，但仍建立完整双幅自由度与真实通道，不使用对称边界；智慧芯位置则应按16个候选索位做离散敏感性，不应把16根\(EA\)\allowbreak{}全部设为相同。随后把测量误差\(\delta y,\allowbreak \delta z,\allowbreak \delta N,\allowbreak \delta m\)\allowbreak{}作为TEST数据加入。没有测量时只能做零均值区间敏感性，不能用目标频率决定差异符号和大小。\par
\phantomsection
\section*{11. 四跨边界、塔顶/辅塔接口与侧跨模态}
\addcontentsline{toc}{section}{11. 四跨边界、塔顶/辅塔接口与侧跨模态}
四跨连续索在北锚、北塔、南塔、辅塔和南锚处的约束直接决定SIDE1—SIDE3和全桥模态排序。至少建立三种独立物理边界工况：\par
B1 固结/不滑移：支点平动按图纸约束，索材料点不跨鞍座滑移，左右跨可具有不同张力； B2 无摩擦滑移：鞍座只约束法向位置，材料可沿索向滑移，鞍座两侧张力连续； B3 摩擦粘滑：粘着状态满足双侧Capstan界限，其中\(\mu_f\)\allowbreak{}为鞍槽摩擦系数，\(\alpha\)\allowbreak{}为包角，\par
\[
e^{-\mu_f\alpha}
\leq\frac{T_2}{T_1}
\leq e^{\mu_f\alpha},
\]
并按静力状态判断粘着或滑移。粘着状态可在线性化后使用对称切线约束；若处于持续滑移或速度相关摩擦，切线一般非保守或非对称，不能直接套用标准实对称特征值问题，需采用相应的非对称/复特征分析或分段状态研究。\par
最终选用哪一档必须由锚固/鞍座详图、施工记录或静载试验决定。由于目标表包含SIDE族，用主跨两端简单固定替代四跨会同时改变边跨局部频率、主跨端部柔度和扭转分支，不能仅凭主跨多数频率接近就认定边界正确。\par
\phantomsection
\section*{12. 全局组装、静力平衡与模态方程}
\addcontentsline{toc}{section}{12. 全局组装、静力平衡与模态方程}
\phantomsection
\subsection*{12.1. 静力切线与广义特征问题}
\addcontentsline{toc}{subsection}{12.1. 静力切线与广义特征问题}
将索、梁、门架、面层、通道和连接件装配后，先求\par
\[
\mathbf R(\mathbf q_0)=\mathbf0.
\]
切线刚度为\par
\[
\mathbf K_T(\mathbf q_0)=
\left.
\frac{\partial\mathbf R}{\partial\mathbf q}
\right|_{\mathbf q_0}
=
\mathbf K_{\mathrm{mat}}+
\mathbf K_{\mathrm{geo}}+
\mathbf K_{\mathrm{conn}}+
\mathbf K_{\mathrm{contact}}.
\]
预应力自由振动方程为\par
\[
\mathbf M(\mathbf q_0)\ddot{\mathbf u}
+
\mathbf K_T(\mathbf q_0)\mathbf u=\mathbf0.
\]
广义特征值问题为\par
\[
\mathbf K_T\boldsymbol\phi_r
=\omega_r^2\mathbf M\boldsymbol\phi_r,
\qquad
\boldsymbol\phi_r^T\mathbf M\boldsymbol\phi_s=\delta_{rs},
\qquad
f_r=\frac{\omega_r}{2\pi}.
\]
建议至少计算前40—60个全局模态，因为十四个目标物理模态未必就是数值排序的前十四阶，边跨、通道和门架局部模态可能插入。由于平移和转角自由度对应的残量单位不同，先用特征长度\(\ell_*\)\allowbreak{}构造对角尺度矩阵\(\mathbf D=\operatorname{diag}(1,\allowbreak 1,\allowbreak 1,\allowbreak \ell_*,\allowbreak \ell_*,\allowbreak \ell_*,\allowbreak \ldots)\)\allowbreak{}，令缩放位移\(\mathbf u^*=\mathbf D\mathbf u\)\allowbreak{}把转角换算为长度量纲，则与之功共轭的缩放残力为\(\mathbf r^*=\mathbf D^{-T}\mathbf r\)\allowbreak{}；每个模态检查\par
\[
r_r^{*}=
\frac{
\|\mathbf D^{-T}(\mathbf K_T\boldsymbol\phi_r-
\omega_r^2\mathbf M\boldsymbol\phi_r)\|_2
}{
\|\mathbf D^{-T}\mathbf K_T\boldsymbol\phi_r\|_2+
\omega_r^2\|\mathbf D^{-T}\mathbf M\boldsymbol\phi_r\|_2
},
\]
并排除单节点局部化、负刚度、松索伪模态和约束残余自由度。\par
\phantomsection
\subsection*{12.2. 不建立全桥有限元的分段连续体动态刚度路线}
\addcontentsline{toc}{subsection}{12.2. 不建立全桥有限元的分段连续体动态刚度路线}
为满足“纯力学、直接计算”的主路线，L1/L2不必依赖全桥有限元网格。每幅猫道在相邻构造站位之间取截面场\par
\[
\mathbf q_p=[v_p,w_p,\theta_p,\chi_{p1},\chi_{p2}]^T,
\]
双幅拼成10维\(\mathbf q=[\mathbf q_U;\allowbreak \mathbf q_D]\)\allowbreak{}，在每个参数均匀或分段常值区间建立\par
\[
\mathbf C\mathbf q''+\omega^2\mathbf M\mathbf q=\mathbf0.
\]
令共轭内力\(\mathbf p=\mathbf C\mathbf q'\)\allowbreak{}，状态\(\mathbf s=[\mathbf q;\allowbreak \mathbf p]\)\allowbreak{}，则长度为\(\ell\)\allowbreak{}的连续段由矩阵指数精确传播：\par
\[
\mathbf s(x+\ell)=
\exp\!\left[
\begin{pmatrix}
\mathbf0&\mathbf C^{-1}\\
-\omega^2\mathbf M&\mathbf0
\end{pmatrix}\ell
\right]\mathbf s(x).
\]
第j个门架、通道或集中质量站位满足\par
\[
\mathbf q_j^+=\mathbf q_j^-,
\qquad
\mathbf p_j^+-\mathbf p_j^-=
(\mathbf K_j-\omega^2\mathbf M_j)\mathbf q_j,
\]
对应跳跃矩阵\par
\[
\mathbf J_j(\omega)=
\begin{bmatrix}
\mathbf I&\mathbf0\\
\mathbf K_j-\omega^2\mathbf M_j&\mathbf I
\end{bmatrix}.
\]
按已恢复的单幅71道门架和全线21道通道位置，把连续段传播矩阵与\(\mathbf J_j\)\allowbreak{}依次相乘，再对北锚、北塔鞍座、南塔鞍座、辅助塔和南锚施加B1/B2/B3边界算子\(\mathbf B(\omega)\)\allowbreak{}，通过\(\det\mathbf B(\omega)=0\)\allowbreak{}或其最小奇异值求根。该方法是解析分段场加有限维矩阵运算，不是全桥有限元；它能在同一框架内执行质量-only、正半定刚度-only、真实通道角度、智慧芯16个候选索位、非对称\(\Delta\mathbf M/\Delta\mathbf K\)\allowbreak{}和SIDE边界消融。正式求根应使用缩放后的最小奇异值、区间计数与连续模态跟踪，避免直接连乘造成的溢出；Galerkin解作为独立交叉校核，而不是共享同一矩阵的重复输出。\par
当前已完成的是L0解析/离散交叉校核、索单元与坐标变换单元测试、16索同波速退化特例、通道质量量级消融和四跨成型线拟合；上述10维分段场、71个门架跳跃矩阵、21个真实通道端口矩阵及B1/B2/B3边界行列式尚未完整组装。因此本文证明了频率尺度与若干机制，但尚未证明TA1、TS1、TS2或严格前14阶数值复现。\par
\phantomsection
\section*{13. 模态分类：先按物理分类，再打开十四项频率}
\addcontentsline{toc}{section}{13. 模态分类：先按物理分类，再打开十四项频率}
\phantomsection
\subsection*{13.1. 能量与质量参与率}
\addcontentsline{toc}{subsection}{13.1. 能量与质量参与率}
对第r阶模态，构件族c的刚度能量比例定义为\par
\[
\eta_{r,c}^{K}=
\frac{\boldsymbol\phi_r^T\mathbf K_c\boldsymbol\phi_r}
{\boldsymbol\phi_r^T\mathbf K_T\boldsymbol\phi_r},
\]
质量比例为\par
\[
\eta_{r,c}^{M}=
\frac{\boldsymbol\phi_r^T\mathbf M_c\boldsymbol\phi_r}
{\boldsymbol\phi_r^T\mathbf M\boldsymbol\phi_r}.
\]
分别统计面层索横向、竖向、截面扭转、畸变、门架、通道、主跨、北侧660 m段、南侧717 m段和最南503 m段能量。L/V/T标签应由主导能量和截面运动定义，SIDE标签应要求主跨外能量占比超过预设阈值并由振型确认，不能按频率邻近命名。\par
\phantomsection
\subsection*{13.2. 主跨纵向奇偶性}
\addcontentsline{toc}{subsection}{13.2. 主跨纵向奇偶性}
整桥不对称，但主跨内部仍可定义关于主跨中点的带符号反射算子\(\mathcal R_m\)\allowbreak{}：极向量和轴向转角在反射下采用各自正确的符号，离散节点先映射到共同观测网格。主跨质量加权奇偶指标为归一化相关系数\par
\[
p_r=
\frac{
\boldsymbol\phi_{r,m}^T\mathbf M_m
\mathcal R_m\boldsymbol\phi_{r,m}
}{
\sqrt{
(\boldsymbol\phi_{r,m}^T\mathbf M_m\boldsymbol\phi_{r,m})
[(\mathcal R_m\boldsymbol\phi_{r,m})^T\mathbf M_m
(\mathcal R_m\boldsymbol\phi_{r,m})]
}
}.
\]
\(p_r\approx+1\)\allowbreak{}表示主跨内近似对称，\(p_r\approx-1\)\allowbreak{}表示近似反对称，中间值表示四跨边界或构造不对称导致混合。S/A标签只用于主跨振型性质，不应把整桥强制成S/A半模型。\par
\phantomsection
\subsection*{13.3. 两幅共同/差动参与率}
\addcontentsline{toc}{subsection}{13.3. 两幅共同/差动参与率}
当上下游两幅具有完全相同质量且通道不含内部质量时，可用\(\mathbf q_+,\allowbreak \mathbf q_-\)\allowbreak{}定义互补参与率。实际不对称质量和多端口通道质量会在共同/差动坐标中产生交叉动能，简单令\(\eta_-=1-\eta_+\)\allowbreak{}会漏项。令\(\mathbf S\)\allowbreak{}表示关于全桥纵向中面\(y\to-y\)\allowbreak{}的带符号反射，而不是无符号的U/D标签互换。极向量分量按极向量变换，转角按轴向量变换；尤其绕顺桥轴的\(\theta_x\)\allowbreak{}在横向镜像下变号。\(\mathbf S\)\allowbreak{}应在共同全局端口观测坐标上构造，并先验证\(\mathbf S^2=\mathbf I\)\allowbreak{}，再用全体系六个刚体平移/转动验证\(\mathbf B_U,\allowbreak \mathbf B_D\)\allowbreak{}与符号映射。令\(\mathbf P_\pm=(\mathbf I\pm\mathbf S)/2\)\allowbreak{}，并用对称参考质量\par
\[
\mathbf M_0=\frac12(\mathbf M+\mathbf S^T\mathbf M\mathbf S)
\]
定义\(E_\pm=(\mathbf P_\pm\boldsymbol\phi)^T\mathbf M_0(\mathbf P_\pm\boldsymbol\phi)\)\allowbreak{}和\(\eta_\pm=E_\pm/(E_++E_-)\)\allowbreak{}。同时另报真实质量下的交叉动能\(E_\times=2(\mathbf P_+\boldsymbol\phi)^T\mathbf M(\mathbf P_-\boldsymbol\phi)\)\allowbreak{}，以及不对称质量和通道内部自由度的未配对动能。这些指标用于区分两幅共同横移、共同竖移、差动竖移/system-roll和单幅局部扭转；横向通道真实柔度使模式往往处于混合状态，不能把所有T模态都解释成两幅刚体差动转角。\par
\phantomsection
\subsection*{13.4. 模态连续跟踪与MAC}
\addcontentsline{toc}{subsection}{13.4. 模态连续跟踪与MAC}
不同模型档次或参数扰动之间先把振型投影到同一组物理观测自由度（相同站位、分量、坐标方向和尺度），再采用参考质量\(\mathbf M_{\mathrm{ref}}\)\allowbreak{}加权MAC\par
\[
\mathrm{MAC}(\boldsymbol\phi_a,\boldsymbol\phi_b)=
\frac{
|\widehat{\boldsymbol\phi}_a^T\mathbf M_{\mathrm{ref}}\widehat{\boldsymbol\phi}_b|^2
}{
(\widehat{\boldsymbol\phi}_a^T\mathbf M_{\mathrm{ref}}\widehat{\boldsymbol\phi}_a)
(\widehat{\boldsymbol\phi}_b^T\mathbf M_{\mathrm{ref}}\widehat{\boldsymbol\phi}_b)
}.
\]
其中帽号表示公共观测投影，不能用一套模型自己的全自由度质量矩阵直接乘另一套不同网格振型。对近重根和避交簇，使用子空间主角度而不强行逐阶跟踪。模式配对成本由MAC、L/V/T能量、主跨奇偶指标、共同/差动参与率和SIDE能量共同组成，频率只作为次要项。附件若没有振型向量，则最终只能依据标签、顺序和频率验收，必须承认频率单独不能唯一证明模型机制正确。\par
\phantomsection
\section*{14. 附件2-3十四项目标的隔离式验收表}
\addcontentsline{toc}{section}{14. 附件2-3十四项目标的隔离式验收表}
以下数据只在模型冻结后加载，来源为附件2-3表4-1：\par
\begingroup
\small
\setlength{\tabcolsep}{3pt}
\setlength{\LTleft}{\fill}
\setlength{\LTright}{\fill}
\begin{longtable}{|>{\raggedright\arraybackslash}p{0.2600\textwidth}|>{\raggedright\arraybackslash}p{0.2600\textwidth}|>{\raggedright\arraybackslash}p{0.2600\textwidth}|}
\hline
\textbf{项目} & \textbf{数值} & \textbf{单位} \\ \hline
\endfirsthead
\hline
\textbf{项目} & \textbf{数值} & \textbf{单位} \\ \hline
\endhead
LS1 & 0.0365 & Hz \\ \hline
VA1 & 0.0700 & Hz \\ \hline
LA1 & 0.0726 & Hz \\ \hline
TA1 & 0.0996 & Hz \\ \hline
VS1 & 0.1028 & Hz \\ \hline
LS2 & 0.1087 & Hz \\ \hline
TS1 & 0.1147 & Hz \\ \hline
SIDE1 & 0.1149 & Hz \\ \hline
SIDE2 & 0.1239 & Hz \\ \hline
VA2 & 0.1438 & Hz \\ \hline
LA2 & 0.1449 & Hz \\ \hline
SIDE3 & 0.1557 & Hz \\ \hline
TS2 & 0.1571 & Hz \\ \hline
VS2 & 0.1744 & Hz \\ \hline
\end{longtable}
\endgroup
验收流程必须是：先保存输入台账和SHA-256；完成静力平衡、网格收敛、质量核对、特征残差和物理分类；生成一个不含目标频率的“预测十四项”文件；再加载上表计算误差。若某一模式没有可靠的物理匹配，应报告“未识别/发生模态重排”，不能用最近频率强行配对。\par
“严格前14阶复现”与“物理家族复现”必须分开发布。严格复现要求冻结模型的全局前14个非刚体根与附件1—14阶同序一一对应，且0.1744 Hz以下没有额外门架、通道、边跨或数值局部根插入；物理家族复现允许在前40—60阶中借助能量、奇偶性、带符号共同/差动指标和MAC找到十四个标签，并允许顺序重排或额外低频模态。后一种结果只能称为“家族级匹配”，不能称为“附件前14阶已复现”。本轮11/14结果是4个L族、4个V族和3个SIDE候选的合成家族级对照；它不是统一L1/L2特征谱，更不构成严格前14阶验收。\par
建议分三级验收。机制级验收要求十四个物理家族、主跨奇偶性、共同/差动性质和SIDE局部性基本正确；工程级初验可暂以十四项平均绝对相对误差不高于5\%、最大误差不高于10\%、T族平均误差不高于7.5\%作为项目目标；加入独立索力、EA、接口和实测几何后，可把目标收紧到平均3\%左右。上述阈值属于本项目建议，不是规范限值，也不能用于反向调参。\par
\phantomsection
\section*{15. 逐级方案与必须记录的扭转诊断量}
\addcontentsline{toc}{section}{15. 逐级方案与必须记录的扭转诊断量}
P0 图纸垂度解析骨架：仅跨长、垂度和重力，输出\(f_1\)\allowbreak{}—\(f_4\)\allowbreak{}及四跨边界敏感性。\par
P1 单幅等效连续索带：用图纸总质量与截面二次矩，检查横/竖/扭转是否出现张力索束退化。\par
P2 单幅16根面层索＋刚性横梁：检验逐索索力差和偏心质量的平移—扭转耦合。\par
P3 将刚性横梁改为图纸大/小横梁，加入面层剪切的四档试验区间并保留畸变坐标，记录T族频率、畸变能量和截面翘曲。\par
P4 加入71道门架及底梁，分别进行门架质量-only、刚度-only、完整门架三组消融；M1/M2按真实站位建立，记录TA1、TS1、TS2及门架能量占比。\par
P5 复制为两幅独立猫道但暂不连接，确认每幅模态成对出现且差异只来自独立几何/索力数据。\par
P6 加入21道通道的端口一致质量-only模型：把局部模型内部质量静力映射/凝聚到真实接口，内部自由度不作为零刚度游离节点保留；观察共同与差动分支的广义质量和频移。\par
P7 加入通过刚体客观性检验的通道端口刚度-only矩阵，保持P6端口自由度集合不变，按真实\(\mathbf B_U,\allowbreak \mathbf B_D\)\allowbreak{}观察共同、差动与扭转分支变化；若新增矩阵正半定，则有序特征值不得下降。这一工况是机制导数，不代表可制造的无质量构件。\par
P8 加入完整Craig—Bampton通道质量和刚度，比较秩一、上下弦多端口、全CB三档的模态连续性。\par
P9 将H1—H12从0°改为图纸实际角度，比较纵向—竖向耦合、塔区模态和T族避交；北侧九道位置采用已恢复的尺寸链坐标，姿态优先由成型线局部切向和端口坐标计算，并与镜像假设对照，不再删除北侧位置。\par
P10 将统一门架节距改为51/57/60 m及塔旁特殊节距，观察局部周期破缺和SIDE族插入。\par
P11 四跨边界从B1切换到B2/B3，重点跟踪SIDE1—SIDE3和主跨端部模态。\par
P12 加入两幅实测索力/几何/附加质量差，检查理想共同/差动模态如何混合。\par
P13 冻结模型并盲比附件十四项；后续任何修订都必须由新增图纸、试验或测量证据触发，并保留修订前后的输入差异、模态MAC与能量迁移。\par
每一步至少输出：总质量及三向质心、关于单幅中心线和双幅中心线的转动惯量、静力索力分布、各支点反力、静力几何残差、前60阶频率、特征残差、L/V/T/畸变能量、门架/通道能量、共同/差动参与率、主跨奇偶指标、四跨能量比例以及与上一步的MAC映射。\par
\phantomsection
\section*{16. 频率偏差的物理诊断规则}
\addcontentsline{toc}{section}{16. 频率偏差的物理诊断规则}
若L/V/T及SIDE全部按近似同一比例偏高或偏低，优先检查单位、总质量、重力、垂度、整体索力和是否重复计质量。若L/V骨架正确而TA1、TS1、TS2整体偏低，优先检查门架横剪/转动端口、通道多端口刚度、面层剪切和连接转动刚度；若T族偏高，检查是否把横梁、网片或通道设成刚性、是否遗漏偏心质量或过度约束两幅。若TA与TS分裂次序错误，检查通道站位、H1—H12姿态、纵向奇偶性和两幅共同/差动耦合。若SIDE族偏差显著而主跨族良好，检查辅塔、边跨长度、鞍座滑移、边跨门架/通道和锚固柔度。若出现几乎成对的完全重根，检查是否无意中施加了两幅完全对称约束；若模式在相近频率间交换，应以MAC和子空间跟踪判断避交，不以阶次编号判断模型“跳错”。\par
对简单、非重根且\(\mathbf K,\allowbreak \mathbf M\)\allowbreak{}均为实对称、参数扰动后边界与自由度集合不变的模态，质量归一化后满足一阶灵敏度\par
\[
\delta\omega_r^2=
\boldsymbol\phi_r^T
\left(\delta\mathbf K-
\omega_r^2\delta\mathbf M\right)
\boldsymbol\phi_r,
\]
从而\par
\[
\frac{\delta f_r}{f_r}
=\frac{1}{2\omega_r^2}
\boldsymbol\phi_r^T
\left(\delta\mathbf K-
\omega_r^2\delta\mathbf M\right)
\boldsymbol\phi_r.
\]
方向判据在纯消融下是严格的。若自由度、边界和\(\mathbf M\)\allowbreak{}保持不变，仅加入客观且正半定的\(\delta\mathbf K\succeq0\)\allowbreak{}，则由Courant—Fischer极值原理，所有按数值排序的特征值都不能下降；若\(\mathbf K\)\allowbreak{}不变，仅加入\(\delta\mathbf M\succeq0\)\allowbreak{}，所有有序特征值都不能上升。“频移方向不能预设”只适用于质量与刚度同时加入、模型替换导致\(\delta\mathbf K\)\allowbreak{}不定、自由度或约束改变，或者按物理标签跨越避交后重新配对的情形。该式用于解释独立证据引起的频移，不用于从频率求\(\delta\mathbf K\)\allowbreak{}或\(\delta\mathbf M\)\allowbreak{}。近重根、精确重根和避交簇必须对投影子空间内的扰动矩阵求特征值，不能逐根套用标量公式；摩擦滑移导致非对称切线时需改用左右特征向量灵敏度。\par
\phantomsection
\section*{17. 网格、数值和独立交叉验证}
\addcontentsline{toc}{section}{17. 网格、数值和独立交叉验证}
索单元长度应先按门架、横梁、通道和支点站位强制分段，再在长区间加密；不得跨越集中质量或连接站位。建议至少做三档网格，使前30个物理模态在中—细网格间频率变化小于0.5\%，T族能量与MAC稳定。梁采用一致质量并保留转动惯量；索与梁质量不能通过节点集中质量再次重复。静力迭代检查力残差、位移增量和总势能；模态检查刚体根数、负特征值、单元能量和特征残差。至少使用两条独立实现交叉核对：一条全三维预应力索—梁FE路径，一条L1/L2能量—Galerkin或独立矩阵组装路径。二者共享图纸输入，但不共享刚度矩阵与求解代码；在简化子问题上先核对单索、16索刚性截面、单门架、单通道和双子结构解析结果，再进入全桥。\par
\phantomsection
\section*{18. 现有材料穷尽检索、补入值与余项处置}
\addcontentsline{toc}{section}{18. 现有材料穷尽检索、补入值与余项处置}
本轮实际检索了用户附件、项目公开归档的51个分卷清单，并从命中分卷提取《猫道图纸1225》124页施工图、《猫道结构复核计算报告0324》118页和《猫道结构抗风性能试验报告》；同时审计站位JSON、构件清单和支承节点表的来源字段。B00、MCT节点/单元/索力/质量/约束及既有ANSYS结果均未进入独立计算。检索结果把原稿多项“待获取”转化为以下已知值、版本冲突或可执行动作。\par
第一，成型线形已经补齐。复核报告表1-9闭合主跨成型垂度227.300 m，表1-8闭合空索垂度218.551 m；总体图255.56 m属于页面上部主缆/桥梁立面，不能作猫道成型垂度。猫道锚点、转索点、塔顶鞍座、两塔下拉点、辅塔和跨中控制点均已取得具体桩号/高程。表1-5跨中桩号K19+836.000的排版笔误由表1-6与表1-9纠正为K18+686.000。\par
第二，全线通道与门架纵向位置已经补齐。21道通道的四跨尺寸链和全部坐标见2.5节；北侧九道不再删去。单幅71道门架的站位也由总体尺寸链闭合。当前真正待核的是北侧九道BOM和局部姿态，不是位置；姿态先由成型索线的局部切向计算，BOM以南侧核心模块5769.2 kg/道和复核报告整道设计恒载10130 kg/道形成两档质量包络。\par
第三，设计材料输入已经补齐。复核报告第6页给猫道承重索\(E=1.20\times10^5\ \mathrm{MPa}\)\allowbreak{}，普通绳金属面积1400.42 mm\(^2\)\allowbreak{}，十六根中一根智慧芯绳面积1292.45 mm\(^2\)\allowbreak{}，十六索束并联\(EA=2.675850\ \mathrm{GN}\)\allowbreak{}；门架六索束名义\(EA=1.0083024\ \mathrm{GN}\)\allowbreak{}。智慧芯绳\(EA\)\allowbreak{}比普通绳低7.71\%，其具体横向安装位置须从M00-04、材料追踪表或安装记录核实；若同应变且位于最外侧\(|y|=2.67\ \mathrm m\)\allowbreak{}，张力形心偏移上限约12.9 mm，归一化平移—扭转耦合约0.7\%，频移可能较小但足以旋转近重根振型。上述值作为CALC-INPUT进入面内理论基线；工作点卸载/再加载切线\(EA_t\)\allowbreak{}仍须由产品证书、张拉—伸长记录或样段试验确认，验证方法与设计基线已明确，不再留空。\par
第四，预应力已有计算初值与独立重算路线。按表1-1均布恒载2.766 kN/m和成型垂度，名义2300 m L0抛物线分布荷载代理给十六索束\(H=8.046711\ \mathrm{MN}\)\allowbreak{}，等分单根0.502919 MN；它不是实际2286.642 m鞍间静力解，且集中质量、连续四跨和下拉系统尚未计入，故该值只是透明代理初值。报告表1-10/1-13已给各跨曲线长度和无应力长度，表1-11/1-14已给恒载各跨索组最大轴力；这些量标记为CALC-AUDIT，而不是继续笼统列为“索力/无应力长度未知”。一项不读取MCT的算术闭合是：主跨猫道索段平均应变\((2321.829-2312.744)/2312.744=0.00392823\)\allowbreak{}，按报告\(E\sum A_i\)\allowbreak{}得索组平均拉力10.5114 MN，较表1-11支点最大11.005 MN低4.49\%；门架主跨段平均应变\((2359.290-2351.835)/2351.835=0.00316987\)\allowbreak{}，按六绳束\(EA\)\allowbreak{}得平均拉力3.1962 MN，较表1-14最大3.348 MN低4.53\%。平均伸长对应平均拉力且应低于支点最大值，方向与量级同时闭合。按同应变，普通/智慧芯单绳平均拉力约660.1/609.2 kN，差50.9 kN。正式静力步骤使用控制点、分布/集中恒载、设计\(EA\)\allowbreak{}、无应力长度和下拉构造重做平衡，不读取MCT初力，也不包含TARGET-SHAPE/TARGET-FREQ。\par
报告后部另有“14628/16=1371 kN”的文字算术错误，正确商数应为914.25 kN；1371 kN不得进入任何单索初力或强度校核。此类报告输出只作CALC-AUDIT，机器复算须保存原式、正确结果和差异标记。\par
第五，门架和通道已经具备从图纸直接建立局部模型的构造路径。复核报告第111页给倒三角通道的上下弦\(\phi152\times6\)\allowbreak{}、竖杆\(\phi102\times4\)\allowbreak{}、斜杆\(\phi51\times4\)\allowbreak{}、Q235、铰接/U形螺栓/抱箍语义；第115页给普通门架梁柱与铰接。由这些几何和独立钢材常数建立三维梁—桁局部模型，先做刚体客观性检查，再生成端口静力凝聚/CB矩阵；U形螺栓、抱箍和网片的半刚性只在局部试验或接触FE后升级，不以十四项频率反标。\par
第六，边界构造已经定位但工作状态需重算。图纸显示锚端球铰拉杆、塔顶连续转索鞍、辅塔连续过鞍和两塔下拉系统，否定“五个平面全部三向固定”的既有抽象。复核报告第9页给两塔下拉点，第105页给4组80 t滑车及基本组合总拉力2200 kN，第109页给44.2°方向；2200 kN不属于已确认的模态恒载状态，故只用于构造/方向核对，工作态预力由施工记录或相应静力组合重算。鞍槽摩擦系数和粘滑状态在现有图纸/报告中没有数值，B1/B2/B3作为有明确上下界含义的三组工况，并以施工记录或静载试验筛选。\par
第七，版本冲突已经隔离。2025.08版《猫道图纸1225》采用6根\(\phi54\)\allowbreak{}门架承重索、829.3 kg门架加313.46 kg底梁；原稿所据2026.05版转录为6根\(\phi50\)\allowbreak{}、1123 kg门架加306.98 kg底梁并出现M1/M2。后续L1/L2主分支拟在取得并175页原PDF并复核后，将较晚版本整套使用；当前只保留带警示的转录诊断分支，旧版作完整对照，任何绳径、门架和底梁都不得跨版本拼装。原稿记录的175页《图纸汇总-2026.08.05.pdf》本体未出现在本轮附件或2026-07-29项目归档中；其SHA-256和转录页码保留但状态为TRANSCRIBED-DRW-A-UNVERIFIED，工程签发前须用原PDF复核。本轮没有因此丢弃已由报告/实体旧图闭合的线形、位置和材料值。\par
第八，测试态差异仍须现场台账确认：两幅实测中心线与索力差、临时电缆/风障/设备/人员、索夹和网片实际状态，以及TARGET-SHAPE振型向量。首轮两幅保持名义相同但自由度完全独立，测试态差异只以有来源的测量更新；没有测量的量进入标注范围的敏感性，不默认为零且不从目标频率倒算。\par
\phantomsection
\section*{19. 数据获取的最小试验包}
\addcontentsline{toc}{section}{19. 数据获取的最小试验包}
钢丝绳试验：每种关键绳型至少一根代表样段，在工作索力的0.6—1.2倍范围进行预循环后测切线\(EA\)\allowbreak{}、滞回和残余伸长；同时保存实际金属面积、绳径和温度。门架试验：选M1/M2各一榀，固定或模拟底部连接，对顶部施加横向力、竖向差动位移和扭矩，识别四端口6自由度柔度矩阵。通道试验/数值：对一整道或一个端部＋中间模块建立局部FE，施加两端共同/差动横移、共同/差动竖移、相对转角和扭矩，生成端口刚度与前若干固定接口模态。面层试验：取包含8—16根索、横梁和网片的代表段，测横向剪切、竖向翘曲和扭转。全桥静态试验：通过已知小质量/小水平力或张拉调整，测两幅猫道的静态位移与索力变化，优先识别边界和连接柔度。所有试验量只用于模型输入或独立验证，不以十四项频率为代价函数。\par
\phantomsection
\section*{20. 数据来源清单}
\addcontentsline{toc}{section}{20. 数据来源清单}
[TRANSCRIBED-DRW-A-UNVERIFIED] 原稿登记《图纸汇总-2026.08.05.pdf》为175页、2026.05版图签，并转录M00-01、M00-04、M01-02、M01-03、M01-05、M04-01、M04-02、M04-03、M05-01—M05-12页码及SHA-256：\nolinkurl{5b6231a766c549db07539df85d8c14b99b816dc0772abb1ab19d7deb782b0bd9}。该原PDF本轮未取得，因此上述信息是待原件复核的转录元数据，不是本轮直接DRW-A证据。\par
[DRW-2025-ARCHIVE] 《00张靖皋长江大桥南航道桥猫道图纸1225.pdf》，124页，MD0-01图签2025.08；从项目发布归档zhangjinggao-full-20260729提取，用于四跨/21道尺寸链交叉核对、锚鞍下拉构造和版本审计。SHA-256：\nolinkurl{8df26c6b04f652b952e6cdf9575f34e76c8d014cbb61edf0b846c1c75d8bc113}。\par
[CALC-REPORT] 《张靖皋长江大桥南航道桥 猫道结构复核计算报告0324.pdf》，118页，2026-03-24；其控制点、材料和荷载表标CALC-INPUT，其成型曲线、索力、无应力长度、位移和反力标CALC-AUDIT。SHA-256：\nolinkurl{8b2f0eb267216cd7d1e9746adb28f8973b4614d27fa1e2290eb12c2a6a45ca90}。\par
[WIND-REPORT] 《附件2-3：猫道结构抗风性能试验报告.pdf》，用于全线21道尺寸链交叉核对、气动节段信息和冻结后的TARGET-FREQ/TARGET-FAMILY验收；其既有ANSYS频率和响应不进入模型输入。SHA-256：\nolinkurl{d17d4061c5726c10b88cc80f3f292b16f0dbf3408e032a30840b1bcaad9173d3}。\par
[ARCHIVE-MANIFEST] Hiram-test/model发布页zhangjinggao-full-20260729及其archive-manifest.csv/package-members.csv；只用于定位上述一手文件和审计派生JSON/CSV的来源边界，不使用归档内B00、MCT或既有求解数值。发布页：\url{https://github.com/Hiram-test/model/releases/tag/zhangjinggao-full-20260729} 。\par
[PROJECT-CONTRACT] 《Catwalk Dynamics 长期工作说明》，用于结果状态、有效性门禁和交付证据规则；本轮读取的实体TXT SHA-256：\nolinkurl{ce1bcf7d7f28db5098653f57b28387a6bdf34ebff4958daa17171b020235d32a}。\par
[TARGET] 附件2-3表4-1，十四项模态标签与频率；仅用于模型冻结后的盲验收，不作为任何输入参数来源。\par
[STD-ROPE] GB/T 20118-2025《钢丝绳通用技术条件》，2025-08-29发布，2026-03-01实施，代替GB/T 20118-2017等；用于钢丝绳产品标记、技术要求、检验、验收和试验口径。该标准不应被误读为给出本项目绳型在工作区间的唯一表观EA。\par
[STD-STEEL] GB 50017-2017《钢结构设计标准》；首轮钢构件线弹性参数可取\(E=2.06\times10^{11}\ \mathrm{Pa}\)\allowbreak{}、\(G=7.90\times10^{10}\ \mathrm{Pa}\)\allowbreak{}、\(\rho=7850\ \mathrm{kg/m^3}\)\allowbreak{}，最终质量以图纸BOM/称重为控制，弹性常数以正式材料标准和材质证明复核。\par
[THEORY-CABLE] Irvine, H. M. and Caughey, T. K., “The linear theory of free vibrations of a suspended cable,” Proceedings of the Royal Society A, 341(1626), 299–315, 1974, DOI: 10.1098/rspa.1974.0189。\par
[THEORY-INCLINED-CABLE] Irvine, H. M., “Free Vibrations of Inclined Cables,” Journal of the Structural Division, ASCE, 104(2), 343–347, 1978, DOI: 10.1061/JSDEAG.0004860。\par
[THEORY-CMS] Craig, R. R. Jr. and Bampton, M. C. C., “Coupling of Substructures for Dynamic Analyses,” AIAA Journal, 6(7), 1313–1319, 1968, DOI: 10.2514/3.4741。\par
[ROPE-MODULUS-GUIDANCE] Bridon技术资料指出钢丝绳具有随绳构造与荷载区间变化的“apparent modulus”，高精度计算宜对实际绳样进行模量试验；该资料只支撑试验必要性，不向本模型提供某个未经本批次验证的数值。\par
\phantomsection
\section*{21. 理论验证执行结果}
\addcontentsline{toc}{section}{21. 理论验证执行结果}
\phantomsection
\subsection*{21.1. 来源纠错对横弯骨架的影响}
\addcontentsline{toc}{subsection}{21.1. 来源纠错对横弯骨架的影响}
独立预测先冻结，随后才加载L族四项目标。结果如下：\par
LS1：抛物线0.036718559 Hz，目标0.0365 Hz，偏差+0.599\%；悬链线0.036768481 Hz，偏差+0.736\%。\par
LA1：抛物线0.073437118 Hz，目标0.0726 Hz，偏差+1.153\%；悬链线0.073114885 Hz，偏差+0.709\%。\par
LS2：抛物线0.110155677 Hz，目标0.1087 Hz，偏差+1.339\%；悬链线0.109560267 Hz，偏差+0.791\%。\par
LA2：抛物线0.146874237 Hz，目标0.1449 Hz，偏差+1.362\%；悬链线0.146028469 Hz，偏差+0.779\%。\par
抛物线四项平均绝对相对误差1.113\%，精确自重悬链线为0.754\%；误用255.56 m主缆垂度时四项平均绝对相对误差为4.641\%。修正垂度由一手控制点触发，目标频率没有参与选择；由同一证据修正后四个L族同时改善，构成来源纠错而非频率拟合的正向验证。\par
\phantomsection
\subsection*{21.2. 面内竖弯验证}
\addcontentsline{toc}{subsection}{21.2. 面内竖弯验证}
5.3节64项Galerkin模型在冻结后对VA1、VS1、VA2、VS2给出0.073437118、0.101932295、0.146874237、0.159158688 Hz，单项误差+4.910\%、−0.844\%、+2.138\%、−8.739\%，平均绝对相对误差4.158\%。该结果通过奇偶性自动分类，没有按最近频率配对。VS1的接近与轴向兼容机制和设计\(EA\)\allowbreak{}量级相容，但单一频率不能单独证明或识别二者；VS2的偏差显示更高阶竖弯已需要集中构件、四跨边界、下拉系统和工作态\(EA_t\)\allowbreak{}，简化模型的适用界限清楚。\par
\phantomsection
\subsection*{21.3. 离散与单元级验证}
\addcontentsline{toc}{subsection}{21.3. 离散与单元级验证}
均匀张弦一致质量有限元相对解析前四阶的最大误差随10、20、40、80、160单元分别为6.632\%、1.652\%、0.412\%、0.103\%、0.0257\%；集中质量从下侧收敛且160单元误差0.0257\%。精确悬链线有限元相对64项Ritz参考的最大误差随20、40、80、160、320单元为1.653\%、0.412\%、0.103\%、0.0257\%、0.00643\%。解析三维桁单元切线与中心差分雅可比的相对Frobenius误差为\(6.94\times10^{-10}\)\allowbreak{}，对称性误差低于\(10^{-14}\)\allowbreak{}；通道局部—全局旋转前后的能量相对误差为\(1.55\times10^{-16}\)\allowbreak{}。八节点链式子结构的Craig—Bampton试验中，保留0个固定接口模态时第一非零特征值误差9.476\%，保留2个降至0.307\%，保留4个降至0.0149\%，全保留时回到机器精度，验证了变换和截断收敛逻辑。\par
\phantomsection
\subsection*{21.4. 四跨成型线拟合与SIDE家族}
\addcontentsline{toc}{subsection}{21.4. 四跨成型线拟合与SIDE家族}
另一路完全不读取目标频率的纯力学核验，直接解析复核报告表1-9的728个成型坐标点，在各开放索段内排除两端1\%塔鞍邻域后做二次曲线拟合，再由\(Hz''=w_x\)\allowbreak{}得到\(H/w_x=1/z''\)\allowbreak{}、\(c=\sqrt{gH/w_x}\)\allowbreak{}及\(f_1=c/(2L_{\mathrm{eff}})\)\allowbreak{}。结果为：北侧开放索段\(L_{\mathrm{eff}}=646.471\ \mathrm m\)\allowbreak{}，\(f_1=0.124225875\ \mathrm{Hz}\)\allowbreak{}，拟合RMS为50.6 mm；主跨开放索段\(L_{\mathrm{eff}}=2286.642\ \mathrm m\)\allowbreak{}，\(f_1=0.036887259\ \mathrm{Hz}\)\allowbreak{}，RMS为77.5 mm；南侧717 m名义跨的开放索段\(L_{\mathrm{eff}}=709.540\ \mathrm m\)\allowbreak{}，\(f_1=0.115652564\ \mathrm{Hz}\)\allowbreak{}，RMS为51.6 mm；最南503 m名义跨的开放索段\(L_{\mathrm{eff}}=496.908\ \mathrm m\)\allowbreak{}，\(f_1=0.166630805\ \mathrm{Hz}\)\allowbreak{}，RMS为38.3 mm。拟合只使用成型几何和重力，未以附件频率选择区间或曲率。\par
冻结后将三个边跨候选按跨域与SIDE标签作家族级对照：南侧717 m段对应SIDE1为0.115652564 Hz，相对0.1149 Hz偏差+0.655\%；北侧段对应SIDE2为0.124225875 Hz，相对0.1239 Hz偏差+0.263\%；最南503 m段对应SIDE3为0.166630805 Hz，相对0.1557 Hz偏差+7.020\%。前两项从独立成型线直接落入目标邻域，第三项的同向偏高则诊断锚固/辅助塔柔度、附加质量或局部边界不能由固定端张弦代理。由于现有附件未提供机器可读振型向量，这三项属于SIDE家族候选，不等于严格阶次配对。\par
采用21.1节精确自重悬链线的4个L族、21.2节浅索Galerkin的4个V族和本节3个SIDE候选，11个非扭转家族的组合平均绝对相对误差为2.508\%，最大误差为VS2的8.739\%。其中SIDE1—SIDE3按南侧717 m段、北侧660 m段、最南503 m段的预声明跨域候选映射对照；在没有机器可读振型向量确认跨域前，2.508\%是该候选映射下的条件统计，不是已确认的11阶统一特征谱误差。这11项来自三个彼此独立、无目标反标的低阶模型；统计用于界定已验证适用域，不能替代尚未组装的统一L1/L2特征谱。\par
\phantomsection
\subsection*{21.5. 静力算术闭合与智慧芯探针}
\addcontentsline{toc}{subsection}{21.5. 静力算术闭合与智慧芯探针}
报告材料、无应力长度和恒载索力之间存在可直接复算的内部闭合。猫道主跨索段平均应变为0.003928234，按十五根普通绳与一根智慧芯绳的设计\(E\sum A_i\)\allowbreak{}得到索组平均拉力10.511365 MN，比表1-11主跨恒载最大11.005 MN低4.486\%；门架主跨索段平均应变为0.003169865，按六绳束\(EA\)\allowbreak{}得平均拉力3.196183 MN，比表1-14最大3.348 MN低4.535\%。平均伸长得到的平均拉力低于支点最大值，方向正确且差值均约4.5\%。同应变下普通绳与智慧芯绳平均拉力约660.141 kN和609.246 kN，相差50.896 kN；这支持把智慧芯具体索位列为16个离散候选位置敏感性，而不是继续把16根\(EA\)\allowbreak{}完全等同。\par
\phantomsection
\subsection*{21.6. 扭转退化的反证}
\addcontentsline{toc}{subsection}{21.6. 扭转退化的反证}
十六根面层索在合成的相同\(H_i/\mu_{x,\allowbreak i}\)\allowbreak{}条件下组装后，前四对平移/扭转特征值在40、80、160单元三档网格的最大相对分裂不超过\(2.39\times10^{-14}\)\allowbreak{}，160单元为数值零，验证7.4节退化定理。该单元测试中的单绳参考水平张力343.432 kN，是人为选取数值等于12.038 kg/m的“名义水平线质量系数”并强制构造\(H_i=\mu_{x,\allowbreak i}c^2\)\allowbreak{}所得；它没有执行实际索长质量到水平投影质量的物理换算，状态为“synthesized\_unit\_test\_only”。该值未分配面层、横梁、网片等附属重力，既不是项目实测索力，也不是复核报告索力。报告同应变普通绳平均拉力约660 kN，按整幅均布恒载与垂度等分的简化水平分力约503 kN/根，二者只作量级边界。\par
本段改用21.4节表1-9成型曲率拟合的Ritz支路作扭转机制基线，其频率与7.4节/主验证JSON中的名义2300 m L0合成退化支路不同；两者均不是统一L1/L2谱，前者是本节发布口径，后者只作几何敏感性。程序先计算并冻结所有候选谐波，然后才打开TARGET-SHAPE。附件图4-5和图4-8可支持TA1、TS1的主跨半波/奇偶判读；附件没有TS2振型图，故\(n=5\)\allowbreak{}是由“第二正对称扭转”的家族名称、奇偶性和顺序推定，并非图中直接观测。采用条件映射\(n=2,\allowbreak 3,\allowbreak 5\)\allowbreak{}后验分类为TA1、TS1、TS2；该操作不是独立预测，也没有参与选基、择阶或参数识别。Ritz支路给出0.073774517、0.110661776、0.184436293 Hz，相对目标0.0996、0.1147、0.1571 Hz的误差为−25.929\%、−3.521\%、+17.401\%。若质量固定、只允许增加同一模态上的刚度，所需残余\(\Delta K/K\)\allowbreak{}分别为+82.266\%、+7.431\%、−27.446\%；TS2要求负刚度增量，与正半定“加构件刚度”矛盾。若刚度固定，仅以质量解释TS2，则需要\(\Delta M/M=+37.829\%\)\allowbreak{}。因此，在固定\(n=2,\allowbreak 3,\allowbreak 5\)\allowbreak{}映射、固定质量与平衡/边界、同模态连续跟踪且不发生避交重标的条件下，三项符号和幅值不一致，严格排除了“对纯张力扭转频率统一乘一个刚度系数”的方案。结论只要求至少引入一项或多项超出统一标量刚度的高阶机制；门架/网片/横梁、真实通道端口、边界、偏心质量、索间波速差及模态重排是待消融候选，并未证明每一项都必需。本节是条件明确的机制否证，不是可接受的三项扭转预测。\par
\phantomsection
\subsection*{21.7. 通道质量消融的适用边界}
\addcontentsline{toc}{subsection}{21.7. 通道质量消融的适用边界}
主跨通道端口一致质量消融给出量级：只加入原稿转录但原PDF未核的H1—H7七道5769.2 kg核心模块时，该分支证据标签为TRANSCRIBED-DRW-A-UNVERIFIED，第一共同插值分支质量比3.037\%，频率降低1.485\%；把主跨十三道同型外推时降低2.568\%；采用复核报告10130 kg/道的设计恒载口径时降低4.382\%。用于分母的342.114 kg/m并非单幅猫道完整总质量，而只是连续索加五类已知离散件均匀化后的“已知部分质量下限”，仍缺面层/底层/扶手网、电缆照明、夹具及其他附件；因此上述质量比和降频绝对值偏大，只能视为上界量级。\par
两节点一致质量插值单元测试中，差动特征向量的附加广义质量为共同向量的三分之一，第一阶降频分别为0.502\%、0.878\%、1.527\%。这个3:1只属于P-A两节点直线连接的质量插值回归，不代表49.65 m三维通道在真实扭转端口上的共同/差动质量比；最终应以Craig—Bampton端口质量矩阵和\(\mathbf B_U,\allowbreak \mathbf B_D\)\allowbreak{}为准。在自由度和刚度不变时，正半定质量-only只能使有序特征值不升；真实通道的完整频移则由质量、端口刚度、内部模态和标签避交共同决定。全21道位置均已纳入站位表，北侧质量/姿态只作为独立场景变化，不再因未核BOM而删除位置。\par
\phantomsection
\subsection*{21.8. 验证结论与发布边界}
\addcontentsline{toc}{subsection}{21.8. 验证结论与发布边界}
已完成并通过或得到可解释误差边界的内容是：输入算术、垂度来源、全21道通道与71道门架站位、L族频率尺度、浅索V族机制、四跨成型线与SIDE候选、索/桁切线、坐标变换、网格收敛、16索退化定理、双幅质量—刚度单元消融和CB降阶。当前状态准确写作“L0与机制单元验证完成，L1/L2尚未求解”。非扭转家族覆盖为11/14、组合平均误差2.508\%；三项T族没有被数值复现，而是用纯张力基线完成了机制否证。统一10维分段连续体、71个门架跳跃矩阵、21个真实通道端口矩阵和B1/B2/B3边界行列式求根仍须完成，故不能签发严格前14阶结果。此处的保留是明确的模型适用域和下一步计算输入清单，不是把已知位置、无应力长度或现有材料留作“未知”。全部机器可读数值、收敛记录和消融结果由随附Python脚本重算并写入JSON；任何MCT/B00初力、节点、约束和频率均未进入计算。按照《Catwalk Dynamics 长期工作说明》的状态定义，本轮属于可审计的本地\(\texttt{partial\_calculation}\)\allowbreak{}，不是已经通过GitHub Actions与完整模型语义门禁的\(\texttt{complete\_supported\_model\_calculation}\)\allowbreak{}。\par
\phantomsection
\section*{22. 最终推荐路线}
\addcontentsline{toc}{section}{22. 最终推荐路线}
最经济且最有判别力的路线为：以本轮已验证的227.300 m成型垂度L/V骨架作为单元回归基准；随后建立L1的单幅16索可畸变连续场，证明张力索束何时会让平移与扭转退化；再按12.2节完整双子结构建立L2，以分段动态刚度/传递矩阵直接传播10维截面场，按图纸尺寸链主分支站位插入71个门架和21个通道跳跃矩阵。门架和通道的端口矩阵优先由图纸杆系的解析梁/桁动态刚度与静力凝聚得到，必要时才用独立局部模型或构件试验交叉校核，不建立全桥有限元。全21道使用已恢复的纵向主分支位置；南侧717 m段采用施工图与复核报告共同支持的178.5+180+180+178.5 m尺寸链，并把抗风报告正文187.5+171+171+187.5 m保留为冲突版本分支；H1—H12使用图纸角度，北侧姿态由成型线切向计算。边界同时按索族保留锚固球铰、连续鞍座、辅助塔和两塔下拉，并对B1/B2/B3粘滑工况做SIDE筛选；在模型冻结前只使用图纸、CALC-INPUT、标准、静力几何和局部试验；最后求取至少60个根、按能量/奇偶/带符号共同差动分类，再打开附件2-3十四项频率。若L/V族接近而T族仍偏差，优先升级门架、通道端口和面层畸变；若SIDE族偏差，优先升级四跨鞍座/锚固/下拉；若所有族同向偏差，再回查版本一致性、完整质量、成型垂度、索力和单位。这个过程把“复现十四项”转化为可审计的物理证据链，避免用频率拟合掩盖模型机制缺失。\par
\phantomsection
\section*{附录A：零污染冻结清单}
\addcontentsline{toc}{section}{附录A：零污染冻结清单}
A1 输入文件中不存在B00/MCT节点、单元、质量、刚度、索力或频率字段。 A2 所有几何与BOM值能回指图纸页码和图号。 A3 钢材与钢丝绳参数能回指正式标准、材料证书或试验报告。 A4 静力平衡只使用重力、安装初应变、支点和独立测量。 A5 TARGET十四项频率文件在求解完成前不可读取。 A6 横向通道质量、门架质量和既有节点集中质量不存在重复计入。 A7 上、下游两幅未施加镜像、等索力或等位移约束。 A8 四跨和辅塔全部存在，SIDE模态不由主跨固定端代理。 A9 门架和通道接口保留转动自由度，凝聚矩阵通过刚体模态、对称性和正半定性检查。 A10 网格收敛、特征残差、模态能量与MAC检查通过后才生成预测表。\par
\phantomsection
\section*{附录B：最核心的三个可证伪假设}
\addcontentsline{toc}{section}{附录B：最核心的三个可证伪假设}
H-T1：若门架与横向通道多端口作用是扭转偏差主因，则在保持质量、自由度、边界及接口映射不变时，由秩一平动连接升级到通过客观性检验且新增量正半定的六自由度端口，应使有序特征值不下降，并使TA1/TS1/TS2出现能由新增端口能量解释的频移或重排，同时L/V骨架变化较小；若按物理标签跟踪后出现下降，必须证明是避交重配而不是把正半定刚度说成“可任意降频”。若T族能量与频率几乎不变，该假设被否证。\par
H-T2：若截面畸变是主因，则由刚性截面升级到至少两个畸变坐标后，T族振型中畸变能量显著增加，频率和模式顺序发生稳定变化；若畸变能量始终很低且频率稳定，该假设被否证。\par
H-T3：若不完美对称和通道逐站姿态是主因，则使用H1—H12实际角度、非均匀门架节距和两幅独立实测状态后，近重根发生可解释分裂或避交；若与全水平、完全对称模型基本相同，则应转向鞍座边界、索EA/索力或局部连接柔度。\par
——文件结束——\par
\end{document}
